% Modernised reading — fonds Alexandre Grothendieck, shelfmark 156-7, pages 1-113.
%
% This is an INTERPRETATION, not a transcription. Everything the critical
% apparatus carried moves into footnotes. In any dispute of substance the
% transcription governs: whoever cites Grothendieck must cite batch-NN.fr.tex.
%
% Pass: Opus 5.5 (claude-opus-5-5), 2026-10-10 — first pass, unchecked against the
% pages by a human. The folder was read whole, its six batches (pages 1-113),
% all mathematics, dated 23 June (p. 1), 24 June (p. 45), 25 June (pp. 54, 69)
% and 26 June 1986 (margin of p. 86, p. 99).
%
% What the folder is. Chapter VII (« GF VII ») of Vers une géométrie des formes,
% « Analysis situs (troisième mouture) ». It follows GF IV (156-4, first
% mouture), GF V (156-5) and GF VI (156-6, second mouture), and the fourth
% mouture (156-8) opens on the « magasins » with which this one ends. Its
% spine is a search for the right primitive data of an « atelier » (algebra of
% figures): (𝔉, ≤, ≪) → (𝓜, ≤, ≪, ≬) → (𝓜, ≤, ≪, ∥ on lieux ℒ) → a dense
% part Λ → « préateliers » → supports and the interior X° of a multistrate →
% the « magasin » (𝓜, ≤, ≪, |∘|), interior disjointness taken as primitive.
% The restarts (pp. 39, 50, 88, 99, 109) are kept as restarts.
%
% Notation fixed for the whole document. 𝔉 the figures, ≤ sub-figure, ≪
% refinement, 𝓜 the multistrates (completely join-irreducible figures), ℒ the
% lieux (≪-minimal non-empty figures); F̃ = {X ∈ 𝓜 | X ≤ F} the strata
% (déploiement); ∂X = X̃ ∖ {X}; ≬ (\between) compatibility, ∥ disjointness,
% ∦ (\nparallel) non-disjointness; Omb X = {Y ∈ 𝓜 | Y ≪ X} (grande ombre),
% omb X = Omb X ∩ ℒ (petite ombre), omb(X)° its interior (lieux not refining a
% proper stratum); Y ≪̊ X interior refinement; F'_L = strata of F' refining L.
% L = X ∩ Y always means the ≤-intersection (common strata). cosupp, supp, ∁,
% Σ_𝓜 for supports (pp. 61 on); |X| ⊂ P a realisation by points (pp. 72 on).
% Three interiors are kept apart: X°_s = supp X ∩ cosupp ∂X (p. 84), X° =
% {Y ∈ Omb X | Y ∥ ∂X} (p. 89), X°_ℓ = Supp omb(X)° (p. 99). |∘| (\mathrel{|{\circ}|})
% interior disjointness (pp. 101 on). The axioms keep his labels, except that
% the second « At pol 1 » (p. 16) and « At spéc » (p. 17) are called At C (pol)
% and At C (spéc) as his own recapitulation of p. 23 calls them.
%
% Substantive departures from the pages, each footnoted where it occurs:
% — p. 5: At 6 is stated for X ≪ F (the page writes X ◁ F, under which it is
%   empty); in the definition of ≪̊ it is Y, not X, that is the smallest element.
% — p. 13: the self-dual form gives that two majorised elements have a SUP, under
%   finiteness of intervals; the page's characterisation of combinatorial simplices
%   (as read) and At simp 1 are necessary, not sufficient (counter-example given).
% — p. 16: L = F ∩ G for the page's F.G.
% — pp. 25-26: Lemme 1 and its corollary need X, Y majorised by a common element
%   (supplied, as Lemme A of p. 40 does); p. 29: the equality omb X ∩ omb Y =
%   omb L comes from ML 1 (page: ML 0), or from (1) and antireflexivity when X ≬ Y.
% — pp. 42, 51, 52, 59: MΛ 1 and its successors stated with X' ≤ X; the vertical
%   sign read ≪ by the transcription would make the axiom fail in polyhedral
%   ateliers; flagged as a doubt.
% — pp. 54-58: the Proposition on φ : 𝔉 → Fig(ℒ) is stated as what the pages
%   actually reach: (i) ⇒ (ii) + φ injective, (ii) ⇒ (i a), (ii) + injectivity
%   ⇒ (i c); (i b) is not reached and the page gives a counter-example.
% — p. 61: D_0 suffices for cosupp A ∈ Σ_𝔉 (the page's « At D » is doubtful).
% — p. 91: X° ∩ Y° = ∅ needs X ≬ Y (as Lemme 4, p. 95, has it).
% — p. 103: Mag 5'' with primes; the derivation of Mag 6 rewritten.
% — p. 107: « équivalence de catégories » read as a correspondence between two
%   kinds of structure on one (𝓜, ≤, ≪); no morphism is defined in the folder.
%
% Modern names given and footnoted as ours: représentation par les éléments
% sup-irréductibles (Birkhoff 1937); structures d'événements premières (Nielsen,
% Plotkin, Winskel 1981) and domaines de configurations; ensembles ordonnés
% simpliciaux (Björner 1984) and complexes simpliciaux abstraits; correspondance
% de Galois / polarité (Birkhoff 1940), espaces d'orthogonalité (Dacey 1968),
% treillis orthocomplémentés, loi orthomodulaire (Husimi 1937); algèbres de
% Boole complètes atomiques; topologie sans points (locales).
%
% What the folder announces and never establishes: the implications b)-d) of
% p. 22 (stated, not proved); the sufficiency of MΛ 0-MΛ 5 (p. 42) and Lemmes
% 2, 3 of p. 49 (« On tourne en rond ! »); (i b) of p. 54; (b) and (c) of the
% Théorème of p. 93 (the abstract lemma meant to give them is false, p. 98);
% the inclusion ⊃ of (a) and statement (b) of p. 100; the converse of the
% Proposition (?) of p. 104; the « magasin polyédral » promised on p. 102.
\documentclass[11pt,a4paper]{article}
\input{../preamble/grothendieck}

\folder{156-7}
\pages{1}{113}
\dating{1986}
\watermark{Édition de démonstration\\interprétation personnelle de l'œuvre}
\foldertitle{[Chapitre] VII. Analysis situs (troisième mouture) : notes manuscrites (23-26/06/1986) — lecture modernisée du dossier entier}

\begin{document}

\section*{Résumé}

\begin{resume}
Prenez un dessin fait de morceaux : des points, des segments, des triangles
collés le long de leurs bords. On peut en tirer un dessin plus petit en
effaçant des morceaux : c'est une \emph{sous-figure}. On peut aussi le
redessiner plus finement, en découpant chaque triangle en petits triangles :
c'est un \emph{raffinement}. Ces cent treize feuillets, écrits du 23 au 26 juin
1986, cherchent à dire ce qu'est une « forme » --- un espace pris du point de
vue de la topologie --- en ne se donnant que ces deux gestes, sans jamais
partir des points de l'espace.

L'objet qu'il manipule s'appelle un \emph{atelier} : un ensemble de figures,
muni des deux relations « est sous-figure de » et « raffine ». Les morceaux
indécomposables d'une figure sont ses \emph{multistrates} ; ceux qu'on ne peut
plus raffiner du tout sont les \emph{lieux}, qui jouent le rôle que les points
jouent d'ordinaire. Deux notions dérivées gouvernent tout : deux figures sont
\emph{compatibles} si elles peuvent vivre ensemble dans une même figure, et
\emph{disjointes} si de plus elles ne partagent aucun morceau. Le chapitre
précédent (dossier 156-6) avait posé une douzaine d'axiomes ; celui-ci les
trie, puis se demande de quel petit nombre de données tout l'atelier se laisse
reconstruire. La réponse change plusieurs fois : d'abord les multistrates et la
compatibilité, puis la seule disjonction entre lieux, puis une partie « dense »
de multistrates, puis une notion de \emph{support} --- ce que la figure
« occupe » ---, enfin une relation nouvelle, « être intérieurement disjoints »,
qui devient la donnée primitive d'une structure qu'il appelle \emph{magasin}.

L'analogie la plus juste est celle des topologies sans points : on y
reconstruit un espace à partir de l'algèbre de ses ouverts, sans se donner
d'abord ses points. Ici l'algèbre est celle des figures, et les points ne
reviennent que comme une réalisation possible parmi d'autres --- les pages 72
à 80 cherchent quand elle existe et quand elle est unique. Une autre analogie
éclaire la notion de support : dans un graphe, l'ensemble des sommets non
voisins d'un ensemble de sommets donné, puis l'ensemble des sommets non
voisins de celui-là, définissent une sorte d'adhérence. Les supports de
Grothendieck sont faits ainsi, à partir de la disjonction, et l'on découvre au
fil des pages qu'ils ne se combinent pas comme les parties d'un ensemble :
« et » ne se distribue pas sur « ou », comme dans les logiques de la mécanique
quantique.

Le dossier est aussi un journal de recherche, et c'est peut-être là son
intérêt le plus large. Un axiome « d'apparence anodine », qui demande que
toute multistrate ait un lieu intérieur, est reconnu comme « l'introduction de
l'infini actuel tout court ! ». Une démonstration échoue deux fois de suite :
« Rien de nouveau ! On tourne en rond ! ». Un lemme général dont il attendait
tout est réfuté par un graphe de quatre sommets : « Rideau… ». Une définition
plausible donne un résultat vide sur l'exemple le plus simple : « C'est triste
mais c'est vrai. » Et la dernière section s'ouvre sur « Je reprends tout, une
nouvelle fois ! ». On voit un auteur qui ne cherche pas à défendre ses
premiers choix, mais à trouver la donnée primitive qui rende les énoncés
naturels --- et qui change de point de vue dès qu'un exemple montre que l'ancien
cache quelque chose.

Ce qu'il bâtit se retrouve aujourd'hui sous plusieurs noms : topologie sans
points, fermetures de Galois et treillis orthocomplémentés, complexes et
ensembles ordonnés simpliciaux, et, de façon inattendue, les structures
d'événements de l'informatique théorique, qui représentent un domaine par des
éléments, un ordre et une relation de compatibilité comme le fait la page 5, à une
condition de finitude près. Le magasin sur lequel le dossier s'arrête est le point de départ de
la quatrième mouture (dossier 156-8).

\keywords{pointless topology, join-irreducible element, prime event structure, simplicial poset, abstract simplicial complex, Galois connection, orthocomplemented lattice, orthomodular law, complete atomic Boolean algebra}
\end{resume}

\section*{Le fil du dossier, et les conventions}

Les feuillets portent sa pagination 1 à 113, qui coïncide avec celle des
archivistes. La page 1 porte « Analysis situs (troisième mouture) » et, dans un
angle, « GF VII ». Le dossier se lit comme une suite de stations ; il se
reprend lui-même plusieurs fois, et ces reprises sont gardées distinctes.

\begin{enumerate}
\item[1.] \emph{Les données d'un atelier et les axiomes du « gros œuvre »}
(pages 1 à 6, 23 juin) : figures, sous-figures, raffinements, multistrates,
lieux ; axiomes At 1 à At 7 ; traduction en termes de multistrates.
\item[2.] \emph{Disjonction et compatibilité} (pages 7 à 12) : ateliers
« spécieux » et « polyédraux ».
\item[3.] \emph{Ensembles ordonnés polyédraux et simpliciaux} (pages 13 et 14).
\item[4.] \emph{Axiomes complémentaires} (pages 14 à 19) : At C, At D, At C
(pol), At C (spéc), At fil, At supp, At ens D.
\item[5.] \emph{Les axiomes des lieux} (pages 19 à 23) : divisibilité At L 1,
At L 2, At L 3, At CL ; tableau des implications, récapitulation.
\item[6.] \emph{Reconstruire un atelier à partir des lieux} (pages 24 à 38) :
axiomes ML, Théorème-Scholie.
\item[7.] \emph{Une partie dense $\Lambda$} (pages 39 à 49, 24 juin à la page
45) : première reprise, qui finit par « On tourne en rond ! ».
\item[8.] \emph{Préateliers} (pages 50 à 53) : récapitulation, trois
axiomatiques.
\item[9.] \emph{Ateliers ensemblistes et quasi-ensemblistes} (pages 54 à 60,
25 juin).
\item[10.] \emph{Supports} (pages 61 à 68).
\item[11.] \emph{Réalisations par des points} (pages 69 à 80, 25 juin à la page
69).
\item[12.] \emph{La ficelle à deux brins} (pages 81 à 83) : les supports ne
sont pas distributifs.
\item[13.] \emph{L'intérieur d'une multistrate} (pages 84 à 100, 26 juin à la
page 99).
\item[14.] \emph{Magasins} (pages 101 à 113) : la relation « intérieurement
disjoints » devient primitive ; « Je reprends tout, une nouvelle fois ! ».
\end{enumerate}

\emph{Les moutures.} Le dossier est la troisième de quatre rédactions d'un même
\emph{Analysis situs} : la première est le chapitre IV (dossier 156-4, 10 juin),
la deuxième le chapitre VI (dossier 156-6, 18-20 juin), la quatrième le
chapitre VIII (dossier 156-8, du 26 juin au 4 juillet). Entre elles, le
chapitre V (dossier 156-5) est une « algèbre des figures ». Celle-ci renvoie
explicitement à la deuxième --- « GF VI, p.~28, 32, 51, 52, 67, 68, 73 » --- ainsi
qu'aux chapitres V et III --- « GF V p.~19, 20 », « GF III ». Elle commence
par récapituler les axiomes de la deuxième, les renumérote, en change le
statut (le recollement des raffinements, axiome At 11 de la deuxième mouture,
devient ici un corollaire), et finit sur une notion, le magasin, dont la
quatrième mouture fait son point de départ.

\emph{Conventions.} $\mathfrak{F}$ est l'ensemble des figures d'un atelier,
$\leq$ la relation « sous-figure de », $\ll$ la relation « raffine ».
$\mathcal{M}$ est l'ensemble des multistrates, $\mathcal{L} \subset \mathcal{M}$
celui des lieux. Pour une figure $F$, $\widetilde{F} = \{X \in \mathcal{M} \mid
X \leq F\}$ est l'ensemble de ses strates, et pour $X \in \mathcal{M}$,
$\partial X = \widetilde{X} \smallsetminus \{X\}$. La compatibilité est notée
$\between$, la disjonction $\parallel$, la non-disjonction $\nparallel$. La
\emph{grande ombre} de $X$ est $\mathrm{Omb}(X) = \{Y \in \mathcal{M} \mid Y \ll
X\}$, sa \emph{petite ombre} $\mathrm{omb}(X) = \mathrm{Omb}(X) \cap
\mathcal{L}$, et $\mathrm{omb}(X)^{\circ}$ l'ensemble des lieux de $X$ qui ne
raffinent aucune strate propre de $X$. Pour deux multistrates $X, Y$, $L = X
\cap Y$ désigne toujours la sous-figure commune, de strates $\widetilde{X} \cap
\widetilde{Y}$ --- jamais l'intersection « fine » $X \cdot Y$. Pour $F' \ll F$ et
$L \leq F$, $F'_{L}$ est la sous-figure de $F'$ formée des strates qui
raffinent $L$.

À partir de la page 84, trois notions d'intérieur d'une multistrate se
succèdent, et un lecteur risque de lire la troisième comme la première. On les
note $X^{\circ}_{\mathrm{s}} = \mathrm{supp}\, X \cap \mathrm{cosupp}\, \partial X$
(page 84), $X^{\circ} = \{Y \in \mathrm{Omb}(X) \mid Y \parallel \partial X\}$
(page 89) et $X^{\circ}_{\ell} = \mathrm{Supp}\, \mathrm{omb}(X)^{\circ}$ (page
99) ; les pages, elles, écrivent $X^{\circ}$ pour les trois. De même le mot
« support » a deux sens : un support \emph{abstrait} est une partie de
$\mathcal{M}$ fermée pour la fermeture de Galois de la disjonction (pages 61 à
68) ; le support \emph{ponctuel} $|X|$ est une partie d'un ensemble de points
$P$ (pages 72 à 80). Enfin la relation « intérieurement disjoints » des pages
101 à 113 est notée $X \mathrel{|{\circ}|} Y$, comme il la dessine.

Les étiquettes des axiomes sont les siennes, avec une exception : deux
énoncés différents portent sur la page le nom « At pol 1 » (pages 13 et 16), et
le second, comme « At spéc » (page 17), est appelé ici At C (pol), resp. At C
(spéc), comme le fait sa propre récapitulation de la page 23. Plusieurs suites
de lemmes recommencent à 1 (pages 21, 25, 48, 55, 89, 97) ; on les cite avec
leur page.

\emph{Ce que le dossier annonce et n'établit pas} : les implications b) à d) du
tableau de la page 22 ; que les axiomes M$\Lambda$ 0 à M$\Lambda$ 5 de la page
42 suffisent ; l'assertion (i b) de la page 54 ; les parties (b) et (c) du
théorème de la page 93 ; l'énoncé (b) et l'inclusion $\supset$ de (a) à la page
100 ; la réciproque de la proposition de la page 104 ; la version « polyédrale »
du magasin annoncée à la page 102.

\pagerange{1}{6}
\section*{Les données d'un atelier et le gros œuvre (pages 1 à 6)}

\subsection*{Les données}

Le 23 juin, il récapitule « les données et axiomes d'une algèbre de figures, ou
atelier », en se limitant aux axiomes purement « algébriques » : ni dimension,
ni connexité, ni régularité topologique, ni combinatoire des simplexes. Parmi
les axiomes, sept sont dits « sine qua non » ; les autres seront « plus ou moins
facultatifs », mais pourront devenir cruciaux « dès qu'il s'agira d'un travail
sur pièces dans le contexte d'un atelier définissant une contrée ». La notion
fondamentale, dit-il, est celle de \emph{figure}, et non celle de multistrate
ou de lieu : on manipule plus commodément les figures si l'on peut traiter
multistrates et lieux comme des figures particulières.

Un atelier est un ensemble $\mathfrak{F}$ muni de deux relations d'ordre,
$\leq$ (« sous-figure ») et $\ll$ (« raffinement »), avec

\[
\textbf{At 0} \qquad F \leq G \Longrightarrow F \ll G ,
\]

axiome « si évident » qu'il ne le numérote pas. Viennent ensuite :
\begin{enumerate}
\item[(a)] la \emph{somme} $F \vee G$, borne supérieure pour $\leq$ quand elle
existe --- c'est alors aussi une borne supérieure pour $\ll$, alors qu'une borne
supérieure pour $\ll$ n'en est pas forcément une pour $\leq$ ;
\item[(b)] l'\emph{intersection} $F \cap G = \mathrm{Inf}^{\leq}(F, G)$, qui
existe toujours, même pour une famille infinie non vide ; quand $F \vee G$
existe, c'est aussi la borne inférieure pour $\ll$ ;
\item[(c)] l'\emph{intersection fine} $F \cdot G = \mathrm{Inf}^{\ll}(F, G)$,
quand elle existe ;
\item[(d)] la \emph{figure vide} $\varnothing_{\mathfrak{F}}$, plus petit élément
pour les deux ordres\footnote{Les quatre lignes sont marquées sur la page de
petits signes cerclés, lus (a) à (d) sans certitude. Les assertions de (a) et
(b) sur $\ll$ ne sont pas démontrées ; elles découlent des axiomes At 4 et At 6
ci-dessous (pour (b) : si $Z \ll X$ et $Z \ll Y$ avec $X, Y$ strates de $F
\vee G$, la plus petite strate de $F \vee G$ que raffine $Z$ est sous $X$ et sous
$Y$). Des notes en marge disent que $F \cdot G$ existe dans le cas ensembliste
$\mathrm{Fig}(\mathcal{L})$ et dans un cas lu « rectiligne ».}.
\end{enumerate}

Les \emph{multistrates} (ou figures élémentaires) sont les figures strictement
irréductibles pour $\leq$\footnote{Nous lisons « strictement irréductible »
comme « complètement sup-irréductible » : $X$ n'est la borne supérieure d'aucune
famille d'éléments $< X$, en particulier $X \neq \varnothing_{\mathfrak{F}}$.
C'est ce que demande l'axiome At 2, qui fait de toute figure la borne supérieure
de ses multistrates.}, avec les ordres induits ; les \emph{lieux} sont les
éléments minimaux de $\mathfrak{F} \smallsetminus \{\varnothing_{\mathfrak{F}}\}$
pour $\ll$, et ce sont des multistrates. Les \emph{strates} d'une figure $F$
sont les multistrates $X \leq F$ ; leur ensemble $\widetilde{F}$ est le
\emph{déploiement} de $F$, et l'on écrit $X \triangleleft F$ (« $X$ incident à
$F$ ») pour $X \in \widetilde{F}$.

Trois lectures de $F$ sont distinguées dès la page 2 : comme partie ordonnée
$\widetilde{F}$ de $\mathcal{M}$ ; comme figure ensembliste dans $\mathcal{M}$,
la \emph{grande multiombre} $\mathrm{Multomb}(F) = \{\mathrm{Omb}(X) \mid X \in
\widetilde{F}\}$, de support la \emph{grande ombre} $\mathrm{Omb}(F)$ ; comme
figure ensembliste dans $\mathcal{L}$, la \emph{petite multiombre}
$\mathrm{multomb}(F)$, de support la \emph{petite ombre} $\mathrm{omb}(F)$.
L'ombre seule ne permet « pratiquement jamais » de reconstituer $F$ ; le
déploiement et la multiombre le permettent.

\subsection*{Les axiomes « sine qua non » (pages 4 et 5)}

Les trois premiers ne portent que sur $(\mathfrak{F}, \leq)$, les quatre suivants
lient $\leq$ et $\ll$\footnote{La numérotation est celle de la récapitulation de
la page 23. Sur la page 5, les chiffres de At 5, At 6, At 7 sont écrits en
surcharge et l'ordre est At 4, At 6, At 5, At 7 ; à la page 6, la traduction en
termes de $\mathcal{M}$ intervertit encore deux numéros par une flèche courbe
($\mathrm{At}_{\mathcal{M}} 6$ est At 5, $\mathrm{At}_{\mathcal{M}} 5$ est At 7).
L'axiome de la plus petite strate était At 5 dans la deuxième mouture (dossier
156-6, p.~21-22).}.

\begin{enumerate}
\item[\textbf{At 1}] $\mathfrak{F} \neq \varnothing$, et toute famille majorée
(pour $\leq$) a une borne supérieure. (Axiome des sommes majorées.)
\item[\textbf{At 2}] Toute figure est la borne supérieure de ses strates : $F =
\mathrm{Sup}_{X \in \widetilde{F}} X$. De façon équivalente, $F \leq G$ si et
seulement si toute strate de $F$ est sous une strate de $G$. (Axiome des
strates.)
\item[\textbf{At 3}] Si toute strate de $F$ est compatible avec toute strate de
$G$, alors $F$ et $G$ sont compatibles --- où $F \between G$ veut dire que $\{F,
G\}$ est majorée, c'est-à-dire que $F \vee G$ existe. (Critère de compatibilité
par strates.)
\item[\textbf{At 4}] $F \ll G$ si et seulement si toute strate de $F$ raffine une
strate de $G$. (Critère de raffinement par strates.)
\item[\textbf{At 5}] Pour des multistrates compatibles $X \between Y$, $X \ll Y
\Rightarrow X \leq Y$. Dans un même déploiement, $\leq$ et $\ll$ coïncident.
\item[\textbf{At 6}] Si $X \in \mathcal{M}$ raffine $F$, l'ensemble $\{Y \in
\widetilde{F} \mid X \ll Y\}$ a un plus petit élément\footnote{La page écrit
l'hypothèse $X \triangleleft F$ ($X$ strate de $F$) ; l'énoncé serait alors vide,
$X$ étant lui-même, par At 5, le plus petit élément de l'ensemble. L'hypothèse
$X \ll F$ est celle de la traduction de la page 6 ($\mathrm{At}_{\mathcal{M}} 4$
a) : « si $X \ll Y$, il y a un minimal dans $\widetilde{Y}^{X}$ ») et du nom que
lui donne la page 23, « axiome de la plus petite strate, ou des raffinements
intérieurs de strates ».}. (Axiome de la plus petite strate.)
\item[\textbf{At 7}] Pour $X, Y, Z \in \mathcal{M}$, $X \mathrel{\overset{\circ}{\ll}}
Y \mathrel{\overset{\circ}{\ll}} Z \Rightarrow X \mathrel{\overset{\circ}{\ll}} Z$.
(Transitivité des raffinements intérieurs.)
\end{enumerate}

Ici $X \mathrel{\overset{\circ}{\ll}} Y$, « $X$ est un raffinement intérieur de
$Y$ », signifie que $Y$ est le plus petit élément de $\{Y' \in \widetilde{Y} \mid
X \ll Y'\}$ : $X \ll Y$, et $X$ ne raffine aucune strate propre de
$Y$\footnote{La page écrit « $X$ est plus petit élément de $\widetilde{Y}^{X}$ » ;
c'est $Y$ qui l'est, comme le dit la glose qui suit (« $X \ll Y' \triangleleft Y
\Rightarrow Y' = Y$ ») et la définition reprise page 25.}. La page 6 ajoute
deux formes de At 6 en termes de multistrates : si $X \ll Y$ il y a un $Z \leq Y$
avec $X \mathrel{\overset{\circ}{\ll}} Z$, et si $X
\mathrel{\overset{\circ}{\ll}} Y$, $X \mathrel{\overset{\circ}{\ll}} Y'$ avec $Y
\between Y'$, alors $Y = Y'$.

\subsection*{Traduction en termes de multistrates (pages 4 à 6)}

At 1 et At 2 disent que $F \mapsto \widetilde{F}$ identifie $\mathfrak{F}$ à un
ensemble $\widetilde{\mathfrak{F}}$ de parties de $\mathcal{M}$ telles que : ce
sont des parties fermées vers le bas pour $\leq$ ; toute partie fermée d'un
élément de $\widetilde{\mathfrak{F}}$ en est un ; elles recouvrent $\mathcal{M}$ ;
la partie vide en est une ; et toute famille majorée a une borne supérieure dans
$(\widetilde{\mathfrak{F}}, \subset)$. Indépendamment de l'ordre, on peut demander
que chaque $X$ appartienne à un plus petit élément de
$\widetilde{\mathfrak{F}}$, qui est alors $\widetilde{X}$.

Quand les figures sont de \emph{type fini} --- réunions finies d'ensembles
$\widetilde{X}$ --- toute la donnée $(\mathfrak{F}, \leq)$ se reconstitue à partir
de $(\mathcal{M}, \leq, \between)$, où la compatibilité est réflexive, symétrique,
et héréditaire :
\[
X \between Y,\ X' \leq X,\ Y' \leq Y \Longrightarrow X' \between Y' .
\]
Les figures sont alors exactement les parties fermées de type fini de
$\mathcal{M}$ formées d'éléments deux à deux compatibles\footnote{C'est, à la
condition de finitude près, la description des \emph{configurations} d'une
\emph{structure d'événements première} (M.~Nielsen, G.~Plotkin, G.~Winskel,
1981) : un ensemble ordonné d'« événements », une relation de conflit symétrique
héritée vers le haut --- ici le contraire de $\between$ ---, et pour
configurations les parties fermées vers le bas sans conflit. At 3 y est la
« cohérence » des domaines de configurations. Rien n'indique que Grothendieck
ait connu ces travaux d'informatique théorique, et les feuillets ne citent
personne. La représentation d'un ordre par ses éléments sup-irréductibles
remonte, pour les treillis distributifs finis, à G.~Birkhoff (1937).}. Dans le
cas général, il faut une donnée de plus : « certains antifiltres infinis »,
c'est-à-dire une notion de famille localement finie de strates. Enfin, avec $\ll$
sur $\mathcal{M}$, l'atelier se lit aussi comme sous-atelier
$\{\mathrm{Multomb}(F)\}$ de l'atelier ensembliste $\mathrm{Fig}(\mathcal{M})$,
par $\mathrm{Omb}(X) = \mathcal{M}_{\ll X}$\footnote{Renvois de la page 6 à la
deuxième mouture : le Scholie de la page 28 et le Théorème-Scholie de la page 32
du dossier 156-6.}.

\pagerange{7}{12}
\section*{Disjonction et compatibilité : ateliers spécieux et polyédraux (pages 7 à 12)}

Il fait passer maintenant les axiomes de disjonction-compatibilité avant ceux
des lieux, « pour ne pas les en faire dépendre ». Deux figures sont
\emph{disjointes} si
\[
F \parallel G \iff F \between G \ \text{et}\ F \cap G = \varnothing_{\mathfrak{F}} ,
\]
notion qui ne dépend que de $\leq$, et qui se teste strate par strate : $F
\parallel G$ si et seulement si $X \parallel Y$ pour toutes strates $X$ de $F$,
$Y$ de $G$. Pour des multistrates, $X \parallel Y$ signifie $X \between Y$ et
qu'aucune multistrate n'est sous les deux.

Le premier axiome tenté (At DC 1, page 7) est barré, et la page 8 en
reformule le contenu. Ce qui reste, sous les noms qu'il fixera à la page 14, est
ceci. Soient $F \between G$, $L = F \cap G$, et des raffinements $F' \ll F$, $G'
\ll G$ ; on note $F'_{L}$ et $G'_{L}$ les sous-figures de $F'$ et $G'$ formées des
strates qui raffinent $L$ --- les raffinements de $L$ « induits » par $F'$ et
$G'$\footnote{La page 8 introduit la lettre $K$ pour $X \cap Y$ et conclut avec
l'indice $L$ ; de même At CD 2 à la page 10. On écrit $L$ partout.}. La
\emph{Proposition 1} souhaitée est
\[
F' \between G' \iff F'_{L} \between G'_{L} ,
\]
dont le sens $\Rightarrow$ est automatique ; le sens $\Leftarrow$ est l'axiome
At CD 1 (futur At C), qui « équivaut à l'ancien At 12 »\footnote{Celui de la
deuxième mouture (dossier 156-6, p.~73-76), où la même équivalence était d'abord
un « Cor.\ Main » puis un axiome. Une note en marge signale que le critère est
« faux dans le cas simplicial », avec un dessin (un losange contenant un cercle
coupé par un diamètre) qui n'est pas repris ici, et qu'il fournit un
contre-exemple à l'axiome At 11 des pages 67-68 de la deuxième mouture.}. Ses
corollaires (Cor.\ Main 1 et 2) : la disjonction est stable par raffinement --- $F
\parallel G$, $F' \ll F$, $G' \ll G \Rightarrow F' \parallel G'$ --- et, pour $F
\between G$, $F' \parallel G' \iff F'_{L} \parallel G'_{L}$.

À la page 9, la forme affaiblie (Propriété $1'_{0}$, axiome At CD $1'_{0}$) ne
demande que la disjonction : sous $F \between G$, $F' \parallel G' \iff F'_{L}
\parallel G'_{L}$, le sens $\Rightarrow$ étant toujours vrai\footnote{L'indice
« $1'_{0}$ » est un petit « o » sous le prime, lecture incertaine.}. Elle donne
qu'un raffinement $F'$ de $F$ avec $F'_{L} = \varnothing_{\mathfrak{F}}$ est
disjoint de tout raffinement de $G$, et qu'en général $F'_{L} =
\varnothing_{\mathfrak{F}}$ équivaut à $F' \parallel L$. « J'ai tendance à
regarder At CD comme axiome obligatoire pour les ateliers. »

Reste le cas où l'on ne suppose pas $F$ et $G$ compatibles. Avec $L = F \cap G$
--- qui existe, la famille des strates communes étant majorée ---, la
\emph{Propriété 2} des ateliers « spécieux » dit : si tout raffinement $F'$ de $F$
avec $F'_{L} = \varnothing$ est disjoint de tout raffinement $G'$ de $G$ avec
$G'_{L} = \varnothing$, alors $F \between G$. Pour les multistrates c'est l'axiome
At CD 2, \emph{critère de compatibilité des ateliers spécieux}. Il n'est pas
satisfait, le plus souvent, par les ateliers \emph{polyédraux} --- ceux où deux
multistrates compatibles qui ont une sous-multistrate commune ont pour
intersection une multistrate --- et il faut pour eux l'axiome At CD 2', qui
demande de plus que $L$ soit vide ou une multistrate. Tous les ateliers qu'il a
en vue satisfont $1'_{0}$ et l'un des deux. La différence est géométrique : dans
le cas spécieux, deux triangles qui ne se rencontrent qu'en deux sommets, sans
partager le côté qui les joint, sont compatibles ; dans le cas polyédral, non.

L'intérêt de ces axiomes, dit la page 10, est de rendre la compatibilité dans
$\mathcal{M}$ superflue : elle se déduit de la disjonction, et l'on peut chercher
les axiomes en termes de $(\mathcal{M}, \leq, \ll, \parallel)$ --- « Attention,
la relation $\parallel$ dans $\mathcal{M}$ ne résulte pas, comme dans
$\mathfrak{F}$, de la seule donnée de $\leq$ ».

\subsection*{Une hésitation sur les simplexes (pages 11 et 12)}

Pour un atelier simplicial, il avait admis que $\ll$ entre simplexes est la
relation ensembliste d'inclusion. Cela fait d'un segment tracé dans un triangle
un raffinement, et accepte comme « subdivision » d'un triangle sa découpe en
deux par un arc joignant deux sommets --- « ce que personne ne regarderait comme
une subdivision simpliciale ! ». Il faut, dit-il, un raffinement plus strict,
comme au chapitre V\footnote{Dossier 156-5, p.~19-20, où le raffinement
« simplicial » $F \overset{s}{\ll} G$ demandait que l'intersection d'une strate
de $F$ et d'une strate de $G$ soit vide ou une sous-strate. Un passage de la page
12 est barré et en partie illisible.}, et une compatibilité « polyédrale »
séparée : l'union de deux figures doit être une figure \emph{polyédrale}. En
résumé, il faut distinguer partout les axiomes des ateliers polyédraux de ceux
des ateliers spécieux, et pour le plongement dans les figures ensemblistes
demander que $\mathrm{Multomb}(F) \overset{\mathrm{pol}}{\ll} \mathrm{Multomb}(G)
\Rightarrow F \ll G$.

\pagerange{13}{14}
\section*{Ensembles ordonnés polyédraux et simpliciaux (pages 13 et 14)}

Un ensemble ordonné $I$ est \emph{polyédral} si deux éléments qui ont un minorant
commun ont une borne inférieure : dès que $I_{\leq x} \cap I_{\leq y} \neq
\varnothing$, cette intersection a un plus grand élément. Quand ces
intersections sont finies (plus généralement inductives), il revient au même de
demander qu'elles soient filtrantes :
\[
x', y' \leq x, y \Longrightarrow \exists\, z,\ x', y' \leq z \leq x, y .
\]
Sous cette forme la condition est autoduale, et si les intervalles
$I_{x' \leq \cdot \leq x}$ sont finis, elle entraîne aussi que deux éléments
majorés ont une borne supérieure\footnote{La page écrit, d'un mot lu avec doute,
« un plus grand élément ». Ce qui est vrai est la borne supérieure, et sous la
finitude des intervalles : si $x, y$ ont deux majorants minimaux $u \neq v$,
ceux-ci ont $x$ pour minorant commun, donc une borne inférieure $w$, qui majore
$x$ et $y$ et est $\leq u$ ; par minimalité $w = u$, et de même $w = v$. Le
majorant minimal est donc unique, et la finitude des intervalles fait que tout
majorant le majore.}. Une figure est polyédrale si son déploiement l'est ---
deux strates non disjointes ont pour intersection une strate --- et un atelier
est polyédral si toutes ses figures le sont ; c'est l'axiome \textbf{At pol 1},
« condition de bonne intersection » : deux multistrates compatibles et non
disjointes ont une borne inférieure dans $(\mathcal{M}, \leq)$.

Un ensemble ordonné est \emph{simplicial} s'il est polyédral et si chaque $I_{\leq
x}$ est un \emph{simplexe combinatoire}, c'est-à-dire isomorphe à l'ensemble
$\mathfrak{P}^{*}(J)$ des parties non vides d'un ensemble fini $J$ ; un ensemble
ordonné fini $S$ avec un plus grand élément est un simplexe combinatoire si et
seulement si $S$ augmenté d'un plus petit élément est un treillis de
Boole\footnote{La page propose une caractérisation dont la fin se lit mal (« ou
que $S$ a un plus grand élément, et que chaque élément non minimal a au moins
deux prédécesseurs »), puis la condition \textbf{At simp 1} : pour toute
multistrate $X$, $\partial X$ est vide ou a au moins deux éléments maximaux.
Ces conditions sont nécessaires, non suffisantes : trois sommets $a, b, c$ et un
élément $m$ au-dessus d'eux, sans arêtes, forment un ensemble polyédral fini à
plus grand élément où $\partial m$ a trois éléments maximaux, et ce n'est pas un
simplexe. La caractérisation donnée dans le texte est la nôtre.}. Une figure,
un atelier sont (combinatoirement) simpliciaux si leurs déploiements le
sont\footnote{Les ensembles ordonnés dont tout intervalle inférieur est un
treillis de Boole sont les \emph{ensembles ordonnés simpliciaux} (A.~Björner,
1984), ensembles de faces des $\Delta$-complexes. La condition polyédrale est
exactement ce qui les ramène aux complexes simpliciaux abstraits : deux arêtes de
mêmes extrémités, ou deux triangles de même bord, ont des minorants communs sans
borne inférieure ; et si les intersections existent, deux éléments de mêmes
sommets sont égaux. Le rapprochement est le nôtre.}.

\pagerange{14}{19}
\section*{Axiomes complémentaires (pages 14 à 19)}

\subsection*{Compatibilité et disjonction des raffinements}

La page 14 donne leurs noms définitifs aux axiomes des pages 8 à 10. Soient $X,
Y, X', Y' \in \mathcal{M}$ avec $X' \ll X$, $Y' \ll Y$, $X \between Y$, $L = X
\cap Y$, et $X'_{L} = \{X'' \in \widetilde{X}' \mid X'' \ll L\}$, $Y'_{L}$ de
même --- l'intersection fine de $X'$ et de $L$\footnote{La page écrit
$\mathrm{Inf}^{\ll}(X, L)$ pour $\mathrm{Inf}^{\ll}(X', L)$.}.

\begin{enumerate}
\item[\textbf{At C}] (critère de compatibilité des raffinements de sous-figures)
$X'_{L} \between Y'_{L} \Rightarrow X' \between Y'$.
\item[\textbf{At D}] (critère de disjonction des raffinements de sous-figures)
$X'_{L} \parallel Y'_{L} \Rightarrow X' \parallel Y'$.
\end{enumerate}

Ils visent les propositions $F' \between G' \iff F'_{L} \between G'_{L}$ et $F'
\parallel G' \iff F'_{L} \parallel G'_{L}$ pour des figures $F \between G$ ---
c'est-à-dire sous-figures d'une même figure $H = F \vee G$. La page 15 les
illustre par le schéma des grandes ombres

\begin{tikzcd}[column sep=small, row sep=small]
 & & \mathrm{Omb}\, H \arrow[dl, no head] \arrow[dr, no head] & & \\
 & \mathrm{Omb}\, F \arrow[dl, no head] \arrow[dr, no head] & & \mathrm{Omb}\, G \arrow[dl, no head] \arrow[dr, no head] & \\
 \mathrm{Omb}\, F' \arrow[dr, no head] & & \mathrm{Omb}\, L \arrow[dl, no head] \arrow[dr, no head] & & \mathrm{Omb}\, G' \arrow[dl, no head] \\
 & \mathrm{Omb}\, F'_{L} & & \mathrm{Omb}\, G'_{L} &
\end{tikzcd}

où chacun des trois carrés est, dit-il, cartésien (une intersection), et le même
pour les déploiements. Il suffit de tester la compatibilité de $F'$ et $G'$ sur
les strates qui raffinent $L$. Le corollaire de At D est \textbf{At
D}$_{0}$ : la disjonction est stable par raffinement\footnote{Si $F \parallel G$,
alors $L = \varnothing_{\mathfrak{F}}$, donc $F'_{L} = G'_{L} = \varnothing$, et At D
donne $F' \parallel G'$.}.

\subsection*{Compatibilité sans compatibilité préalable}

Sans supposer $X \between Y$, on forme encore $L = X \cap Y$. Les deux critères
des pages 9 et 10 deviennent :
\begin{enumerate}
\item[\textbf{At C (pol)}] si tout raffinement $X'$ de $X$ avec $X'_{L} =
\varnothing$ est disjoint de tout raffinement $Y'$ de $Y$ avec $Y'_{L} =
\varnothing$, et si de plus $L$ est vide ou une multistrate, alors $X \between Y$ ;
\item[\textbf{At C (spéc)}] le même sans la condition sur $L$\footnote{Sur la page
le premier s'appelle « At pol 1 », numéro écrit en surcharge, nom déjà pris par
la condition de bonne intersection de la page 13 ; le second « At spéc ». Une
longue note en marge, lue en partie, commence par « On ne veut pas cet
axiome ».}.
\end{enumerate}
Sous At C (pol), deux figures $F, G$ sont compatibles si et seulement si (i)
deux strates de $F$ et de $G$ qui ont un minorant commun ont pour intersection
une strate, et (ii) les raffinements de $F$ et de $G$ qui évitent $L = F \cap G$
sont disjoints\footnote{La page écrit $L = F.G$, avec le point de
l'intersection fine ; les axiomes sur lesquels repose l'énoncé emploient $F \cap
G$.} ; sous At C (spéc), la seule condition (ii). Et l'on a
\[
\text{At D} + \text{At C (spéc)} \ \text{ou}\ \text{At D} + \text{At C (pol)}
\Longrightarrow \text{At C} \Longrightarrow \text{At D} \Longrightarrow
\text{At D}_{0} ,
\]
où, dit une note en marge, la forme affaiblie de At D qui ne traite que le cas
$X'_{L} = Y'_{L} = \varnothing$ suffirait.

\subsection*{Sommes filtrantes, supports, cas ensembliste}

\textbf{At fil} (critère des sommes filtrantes localement finies) : une famille
filtrante croissante $(F_{i})$ de raffinements de $F$, localement finie dans $F$
--- pour chaque strate $X$ de $F$, la famille des $F_{i}|X$ est stationnaire ---
a une borne supérieure. Moyennant At C, il en tire le \emph{recollement des
raffinements}, en deux versions selon qu'on recolle sur des sous-figures ou sur
des strates\footnote{C'était l'axiome At 11 de la deuxième mouture (dossier
156-6, p.~67-70) ; il devient ici un corollaire. Le Corollaire 2 de la page 18
énonce l'analogue de At fil pour des familles croissantes à la fois pour $\ll$ et
pour $\leq$. Une note en marge se demande
s'il faut supposer les $F_{i}$ croissants pour $\ll$ aussi, « et prendre Sup
pour $\ll$ ?? ».}.

\textbf{At supp} (axiome des supports d'un raffinement) : si $G \ll X$ et qu'il
existe une multistrate $Y$ disjointe de $G$ et non disjointe de $X$, on peut en
choisir une qui raffine $X$. Intuitivement : si $G$ ne remplit pas $X$, un
morceau de $X$ en est témoin. La Proposition qui suit dit le même pour une figure $F$
quelconque au lieu de $X$. La question qui suit --- peut-on trouver $Y' \ll Y$
avec $Y' \ll F$ ? --- se ramène, pour un lieu $x$, à $x \nparallel Y \iff x \ll
Y$, ce qui « n'est guère vérifié que pour les ateliers ensemblistes ». D'où
\textbf{At ens D} : $X \parallel Y$ si et seulement si $\mathrm{Omb}(X) \cap
\mathrm{Omb}(Y) = \varnothing$, avec pour corollaires $F \parallel G \iff
\mathrm{Omb}\, F \cap \mathrm{Omb}\, G = \varnothing$, $x \parallel F \iff x
\notin \mathrm{omb}(F)$, et, pour deux lieux, $x \parallel y \iff x \neq y$. At
ens D entraîne At supp ; « At supp est vérifié dans les ateliers auxquels j'ai
songé, mais non At ens D ».

\pagerange{19}{23}
\section*{Les axiomes des lieux (pages 19 à 23)}

\textbf{At L 1} (axiome de divisibilité) : toute multistrate a un lieu
intérieur, $\mathrm{omb}(X)^{\circ} \neq \varnothing$, c'est-à-dire qu'il existe
$x \in \mathcal{L}$ avec $x \mathrel{\overset{\circ}{\ll}} X$\footnote{C'était
l'axiome des lieux At 8 de la deuxième mouture, formule (1.52) (dossier 156-6,
p.~45).}. Cet axiome d'apparence anodine, dit-il, implique que les espaces
qu'on décrit par des ateliers sont « infiniment divisibles », et donc infinis
quand ils ne sont pas discrets : « c'est l'introduction de l'infini actuel tout
court ! »\footnote{L'infinitude ne suit pas de At L 1 seul : il faut encore que
les multistrates soient assez nombreuses, que chaque lieu intérieur découpe $X$
en multistrates de l'atelier. L'exemple de la page 69 le confirme dans le cas
d'un ensemble ordonné : At L 1 y équivaut à la densité de l'ordre, et un ordre
dense à deux éléments est infini.}

\textbf{At L 2} (critère de disjonction par les lieux) : $X \parallel Y$ si et
seulement si $x \parallel y$ pour tous $x \in \mathrm{omb}(X)$, $y \in
\mathrm{omb}(Y)$ --- « à ne pas confondre avec $\mathrm{omb}(X) \cap
\mathrm{omb}(Y) = \varnothing$ ! »\footnote{C'était At 9 dans la deuxième
mouture (dossier 156-6, p.~61). La mise en garde compte : dans un atelier non
ensembliste deux lieux distincts peuvent ne pas être disjoints.} Il vaut alors
pour les figures.

\textbf{At L 3} (critère de disjonction pour les lieux) : si $X \between Y$ et $X
\neq Y$, deux lieux intérieurs $x \in \mathrm{omb}(X)^{\circ}$, $y \in
\mathrm{omb}(Y)^{\circ}$ sont disjoints. Autrement dit, dans une figure, des lieux
intérieurs à deux strates distinctes sont disjoints ; moyennant At D, il suffit
de le demander pour $Y \leq X$.

\textbf{At CL (spéc)} (critère de compatibilité par les lieux, ateliers
spécieux) : si $x \parallel y$ pour tous $x \in \mathrm{omb}(X) \smallsetminus
\mathrm{omb}(L)$, $y \in \mathrm{omb}(Y) \smallsetminus \mathrm{omb}(L)$, où $L = X
\cap Y$, alors $X \between Y$. Une note en marge se demande s'il vaut mieux poser
« il faut et il suffit ».

La démonstration, à cheval sur les pages 20 et 21, montre que \emph{sous At L 2,
At C (spéc) équivaut à At CL (spéc)}. La condition sur les lieux
\[
(\mathrm{AL}) \quad \forall\, x \in \mathrm{omb}(X) \smallsetminus
\mathrm{omb}(L),\ y \in \mathrm{omb}(Y) \smallsetminus \mathrm{omb}(L) :\ x
\parallel y
\]
équivaut à la condition sur les raffinements
\[
(\mathrm{A}) \quad \forall\, X' \ll X,\ Y' \ll Y \ \text{avec}\ X'_{L} = Y'_{L} =
\varnothing :\ X' \parallel Y' .
\]
\footnote{La page écrit, dans (A), « $X \parallel Y$ » pour $X' \parallel Y'$, et
pour l'implication inverse « en prenant $X, Y \in \mathcal{L}$ » pour $X', Y' \in
\mathcal{L}$.}
De (AL) à (A) : par At L 2 il suffit de voir $x \parallel y$ pour $x \in
\mathrm{omb}(X')$, $y \in \mathrm{omb}(Y')$, et $X'_{L} = \varnothing$ force $x
\notin \mathrm{omb}(L)$. De (A) à (AL) : il suffit de prendre $X', Y'$ des lieux.
On a utilisé deux lemmes :
\begin{enumerate}
\item[] \emph{Lemme 1} (page 21). Pour $P \leq Q$ et $x \in \mathrm{omb}(Q)$, $x_{P}
= \varnothing$ si et seulement si $x \notin \mathrm{omb}(P)$ --- car un lieu
n'a pas d'autre strate que lui-même, donc $x_{P}$ vaut $x$ ou $\varnothing$.
\item[] \emph{Lemme 2} (page 21). Pour $P \leq Q$ et $F \ll Q$, $F_{P} = \varnothing$
si et seulement si $F \parallel P$ --- le sens $\Rightarrow$ par At D, l'autre
avec les seuls At 1 à At 7\footnote{En effet $F \parallel P$ entraîne $F_{P}
\parallel P$ ; comme $F_{P} \ll P$ et $F_{P} \between P$, At 5 donne $F_{P} \leq P$,
donc $F_{P} = F_{P} \cap P = \varnothing$.}.
\end{enumerate}
La version polyédrale \textbf{At CL (pol)} ajoute la condition que $L$ soit vide
ou une multistrate.

\subsection*{Le tableau des implications et la récapitulation (pages 22 et 23)}

La page 22 enregistre :
\begin{enumerate}
\item[a)] At L 2 + At CL (spéc ou pol) $\iff$ At L 2 + At C (spéc ou pol) ;
\item[b)] At D$_{0}$ + At L 1 + At CL $\Rightarrow$ At L 2 ;
\item[c)] At L 2 + At L 3 $\Rightarrow$ At D $\Rightarrow$ At D$_{0}$ ;
\item[d)] At L 3 + At CL (spéc) $\Rightarrow$ At C ;
\item[e)] At D $\Rightarrow$ At L 3\footnote{a) est ce que démontrent les pages
20-21. e) se vérifie aussitôt : si $x \in \mathrm{omb}(X)^{\circ}$ et $y \in
\mathrm{omb}(Y)^{\circ}$ avec $X \neq Y$ compatibles, l'un au moins de $x_{L}$,
$y_{L}$ est vide et l'autre est vide ou égal au lieu lui-même, et At D conclut.
c) se vérifie en plaçant chaque lieu dans sa plus petite strate (At 6) et en
appliquant At L 3 ou, si les deux strates coïncident, l'hypothèse sur $X'_{L}$,
$Y'_{L}$ --- pourvu que $X'_{L}$ soit bien la borne inférieure de $X'$ et de $L$
pour $\ll$, comme le veut la page 14. b) et d) ne sont qu'énoncées.}.
\end{enumerate}
« Pratiquement », retient-il, $C \Rightarrow D \Rightarrow D_{0}$, et L 2 et L 3
entraînent $D$ ; pour obtenir $C$ à partir des lieux il faut At CL. Ainsi
\[
\text{At L 2} + \text{At L 3} + \text{At CL (spéc)}
\]
implique tous les axiomes de compatibilité-disjonction de style spécieux, et avec
At L 1 on a tout. La page 23 récapitule l'ensemble : At 1 à At 7, axiomes du
« gros œuvre » ; At C, At D, At D$_{0}$, At C (spéc), At C (pol), axiomes de
compatibilité et de disjonction ; At Fil et At supp ; At L 1 à At L 3, At CL
(spéc), At CL (pol), axiomes concernant les lieux.

\pagerange{24}{38}
\section*{Reconstruire un atelier à partir des lieux (pages 24 à 38)}

Le critère At CL réduit la connaissance de la compatibilité dans $\mathcal{M}$ à
celle de la disjonction entre lieux, c'est-à-dire à une structure de graphe sur
$\mathcal{L}$ --- dans un atelier ensembliste, rappelle-t-il, $x \parallel y \iff
x \neq y$. D'où la question : partant de $(\mathcal{M}, \leq, \ll)$, de
$\mathcal{L}$ l'ensemble des éléments minimaux pour $\ll$, et d'une relation
symétrique antiréflexive $\parallel$ sur $\mathcal{L}$, à quelles conditions
proviennent-ils d'un atelier spécieux satisfaisant At L 1 à At L 3 et At CL
(spéc) ? La compatibilité sera \emph{définie} par
\[
(1) \qquad X \between Y \iff \forall\, x \in \mathrm{omb}(X) \smallsetminus
\mathrm{omb}(L),\ y \in \mathrm{omb}(Y) \smallsetminus \mathrm{omb}(L) :\ x
\parallel y ,
\]
où $\mathrm{omb}(L) = \bigcup_{Z \leq X, Y} \mathrm{omb}(Z)$.

\subsection*{Les axiomes ML (pages 25 à 34)}

On pose, pour $X \in \mathcal{M}$, $\mathrm{omb}(X) = \{x \in \mathcal{L} \mid x
\ll X\}$ et $\mathrm{omb}(X)^{\circ} = \mathrm{omb}(X) \smallsetminus \bigcup_{Y
< X} \mathrm{omb}(Y)$, et $Y \mathrel{\overset{\circ}{\ll}} X$ comme plus haut.
Dans la numérotation définitive de la page 34 :
\begin{enumerate}
\item[\textbf{ML 0}] $X \leq Y \Rightarrow X \ll Y$ (« axiome trivial ») ;
\item[\textbf{ML 1}] (raffinements induits) si $Z \ll X$, $X' \leq X$, et si un
lieu $x$ raffine $Z$ et $X'$, il existe $Z' \leq Z$ avec $Z' \ll X'$ et $x \ll
Z'$ ; si de plus $Z \leq X$, on peut prendre $Z' \leq X'$ ;
\item[\textbf{ML 2}] (divisibilité) $\mathrm{omb}(X)^{\circ} \neq \varnothing$
pour tout $X$ ;
\item[\textbf{ML 3}] (plus petite strate) si $x \ll X$, il existe $X' \leq X$ avec
$x \mathrel{\overset{\circ}{\ll}} X'$ ;
\item[\textbf{ML 4}] $\parallel$ est symétrique et antiréflexive ;
\item[\textbf{ML 5}] (disjonction des lieux) si $Y \leq X$, $y \in
\mathrm{omb}(Y)$ et $x \in \mathrm{omb}(X) \smallsetminus \mathrm{omb}(Y)$, alors
$x \parallel y$\footnote{Le dossier numérote trois fois ces axiomes. La page 25
pose ML 0 et ML 1 (sous la forme $X' \leq X \geq X''$), la page 26 ML 2 et ML 3,
la page 29 appelle « ML 0 » la symétrie de $\parallel$ et la page 30 « ML 4 » la
disjonction des lieux ; la page 34 renumérote ML 0 à ML 5, et une note de bas
de page y dit « l'axiome ML 4 est vide » quand $\leq$ est discret --- c'est ML 5,
l'ancien ML 4, qui l'est.}.
\end{enumerate}
ML 2 et ML 3 disent que les $\mathrm{omb}(Y)^{\circ}$, $Y \leq X$, forment une
partition de $\mathrm{omb}(X)$ ; ML 2 est At L 1, ML 3 la forme de At 6 pour les
lieux, ML 5 la forme de At L 3. Une note dit que ML 2 entraîne que si $\leq$ et
$\ll$ sont égales, $\mathcal{M}$ est discret.

La chaîne de lemmes est la suivante.
\begin{enumerate}
\item[] \emph{Lemme 1} (page 25). Pour $X, Y$ majorés par un même
élément\footnote{Hypothèse ajoutée : la démonstration applique ML 1 à $X$ et $Y$,
qui le demande, et sans elle l'équivalence de (i) et (ii) ferait coïncider $\ll$ et
$\leq$. La reprise de la page 40 (Lemme A) la pose explicitement, et le Lemme 4
de la page 30 la remplace par $X \between Y$.}, sont équivalents : $X \leq Y$ ; $X
\ll Y$ ; $\mathrm{omb}\, X \subset \mathrm{omb}\, Y$ ;
$\mathrm{omb}(X)^{\circ} \subset \mathrm{omb}(Y)$ ;
$\mathrm{omb}(X)^{\circ} \cap \mathrm{omb}(Y) \neq \varnothing$. En effet, si
$x$ est dans la dernière intersection, ML 1 donne $Z \leq X, Y$ avec $x \ll Z$,
et $x \mathrel{\overset{\circ}{\ll}} X$ force $Z = X$. \emph{Corollaire} : sous la
même hypothèse, deux multistrates distinctes ont des intérieurs disjoints, donc
la strate $Y$ de ML 3 est unique et c'est le plus petit élément de
$\{X' \leq X \mid x \ll X'\}$.
\item[] \emph{Lemme 2} (page 27). Admettant ML 1 sous la forme généralisée
ci-dessus : si $Z \ll X \geq X'$ et $x \mathrel{\overset{\circ}{\ll}} Z$ avec $x \ll
X'$, alors $Z \ll X'$, et $x \mathrel{\overset{\circ}{\ll}} X' \iff Z
\mathrel{\overset{\circ}{\ll}} X'$.
\item[] \emph{Corollaires 1 à 4} (pages 28 et 29). Tout raffinement $Z \ll X$ est
raffinement intérieur d'une unique strate de $X$ ; $Z \ll X'$ se teste par
$\mathrm{omb}(Z)^{\circ} \cap \mathrm{omb}(X') \neq \varnothing$, et $Z
\mathrel{\overset{\circ}{\ll}} X'$ par $\mathrm{omb}(Z)^{\circ} \subset
\mathrm{omb}(X')^{\circ}$ ; d'où la transitivité des raffinements intérieurs
(Corollaire 4), c'est-à-dire At 7.
\end{enumerate}
Il fait ensuite intervenir $\parallel$ et la définition (1). Pour $X, Y$ majorés par
un même élément, ML 1 donne $\mathrm{omb}(X) \cap \mathrm{omb}(Y) =
\mathrm{omb}(L)$\footnote{La page invoque « ML 0 » ; c'est ML 1 (page 25) qui
donne cette égalité. Pour $X \between Y$ quelconques elle vient de (1) et de
l'antiréflexivité : un lieu commun hors de $\mathrm{omb}(L)$ serait disjoint de
lui-même. C'est ainsi qu'elle est employée dans le Lemme 4.}, d'où le
\emph{Lemme 3} (deux sous-multistrates d'une même multistrate sont compatibles,
« sera contenu dans lemme 5 »), le \emph{Lemme 4} (sous $X \between Y$, le Lemme 1
vaut), le \emph{Lemme 5} (la compatibilité définie par (1) est héréditaire pour
$\leq$, démonstration par ML 1 et ML 5), le \emph{Lemme 6} et son corollaire
(si $X \between Y$, $Z \mathrel{\overset{\circ}{\ll}} X$ et $Z \ll Y$, alors $X \leq
Y$ ; en particulier, si $Z$ est raffinement intérieur de deux multistrates
compatibles, elles sont égales, c'est $\mathrm{At}_{\mathcal{M}} 4$ b), le \emph{Lemme 7} (pour deux lieux,
$x \between y$ si et seulement si $x = y$ ou $x \parallel y$, et $\parallel$ est la
disjonction au sens de l'atelier construit) et le \emph{Lemme 8} (At L 2,
« tautologique »).

\subsection*{Le Théorème-Scholie (pages 33 à 37)}

\textbf{Théorème-Scholie.} Soient $(\mathcal{M}, \leq, \ll)$ un ensemble muni de
deux relations d'ordre, $\mathcal{L}$ l'ensemble de ses éléments minimaux pour
$\ll$, et $\parallel$ une relation symétrique antiréflexive sur $\mathcal{L}$.
Ces données proviennent d'un atelier spécieux satisfaisant At 1 à At 7, At L 1 à
At L 3 et At CL (spéc) si et seulement si ML 0 à ML 5 sont satisfaits.

L'atelier construit est l'ensemble $\underline{\Phi}$ des parties de $\mathcal{M}$
fermées pour $\leq$ et formées d'éléments deux à deux compatibles au sens de (1),
avec l'inclusion pour $\leq$ et $F \ll G$ si tout élément de $F$ raffine un
élément de $G$ ; les $\widetilde{X}$ en sont les figures élémentaires. Tout
atelier qui correspond à ces données est un sous-atelier spécieux de
$\underline{\Phi}$ : une partie $\mathfrak{F}$ stable par passage aux parties
fermées, par réunion de deux figures compatibles, et qui contient les
$\widetilde{X}$\footnote{La page écrit $\{\{X\}\}$ ; d'après la page 35, la figure
élémentaire attendue est $\widetilde{X}$.} --- donc toutes les parties fermées des
réunions finies de $\widetilde{X}$. « À part un petit flottement dû à la
possibilité de familles infinies », tout s'exprime donc par $(\mathcal{M}, \leq,
\ll, \parallel)$ avec $\parallel$ défini sur $\mathcal{L}$ seulement, soumis à six
axiomes, « au lieu de la douzaine d'axiomes At 1-7, At L 1-3, At CL
(spéc) »\footnote{La page 36 parle de « quatre axiomes (cinq avec (ML 0)) ML 1
-- ML 4 », d'après l'ancienne numérotation ; la page 34 en énumère six.}.

Si l'on se donne $\parallel$ sur $\mathcal{M}$ et non sur $\mathcal{L}$, il faut un
axiome de plus, ML 6 : $X \parallel Y$ si et seulement si tous les lieux de $X$
sont disjoints de tous ceux de $Y$. La variante \emph{polyédrale} change un peu la
définition (1) et demande ML 0 pol : chaque $\widetilde{X}$ est un ensemble ordonné
polyédral au sens de la page 13.

\subsection*{La marge de choix (page 38)}

Partant de $(\mathcal{M}, \leq, \ll)$, le choix de $\parallel$ est très libre. Si
$\Gamma_{0}$ est l'ensemble des paires $\{x, y\}$ de lieux pour lesquelles il
existe $Y \leq X$ avec $y \ll Y$, $x \ll X$, $x \not\ll Y$ --- les paires que ML 5
force à être disjointes ---, toute structure de graphe $\Gamma$ avec $\Gamma_{0}
\subset \Gamma \subset \mathfrak{P}_{2}(\mathcal{L})$ convient, puisque seuls ML 4
et ML 5 parlent de $\parallel$. « Ça fait de la marge ! » En revanche, poser $x
\parallel y \iff x \neq y$ ne suffit sans doute pas à rendre l'atelier
ensembliste : ML 0 à ML 3 n'impliquent ni que $\mathcal{M} \to
\mathrm{Fig}(\mathcal{L})$ soit injective, ni que $\ll$ soit induite par la
relation ensembliste.

\pagerange{39}{49}
\section*{Une partie dense (pages 39 à 49)}

\subsection*{Première reprise (pages 39 à 44)}

La minimalité des lieux n'a servi qu'à la fin, observe-t-il, pour montrer que
la relation de départ est la disjonction de l'atelier construit. « On va tout
reprendre. » Une partie $\Lambda \subset \mathcal{M}$ est \emph{dense} si toute
multistrate a un raffiné intérieur dans $\Lambda$ ; $\Lambda = \mathcal{M}$ est
l'exemple « universel », $\Lambda = \mathcal{L}$ l'est sous At L 1. On pose
$\mathrm{omb}_{\Lambda}(X) = \{Z \in \Lambda \mid Z \ll X\}$ et
$\mathrm{omb}_{\Lambda}(X)^{\circ}$ les éléments de $\Lambda$ qui sont raffinés
intérieurs de $X$. Les axiomes M$\Lambda$ 0 à M$\Lambda$ 3 sont ML 0 à ML 3 avec
$\Lambda$ au lieu de $\mathcal{L}$, et les Lemmes A, B, C et leur corollaire
(page 40) redisent les lemmes et corollaires des pages 25 et 28, le Lemme A avec
l'hypothèse explicite d'un majorant commun\footnote{Dans le corollaire de la page
40, la condition (iv) garde l'écriture « $Z \ll X'$ » de la page 28 pour $Y \ll
X'$.}.

La compatibilité est maintenant donnée sur $\mathcal{M}$, et la disjonction
définie par $X \parallel Y \iff X \between Y$ et $\widetilde{X} \cap \widetilde{Y}
= \varnothing$. Le raisonnement de la page 30 a besoin de $\mathrm{omb}_{\Lambda}(X)
\cap \mathrm{omb}_{\Lambda}(Y) = \mathrm{omb}_{\Lambda}(L)$, qu'il faut poser en
axiome (Lemme A', page 41) ; il en tire, comme pages 32 et 33, que si un même
élément est raffinement intérieur de deux multistrates compatibles, celles-ci sont
égales (Lemme B et son corollaire).
La \emph{Proposition} de la page 42 énonce les propriétés
caractéristiques d'un système $(\mathcal{M}, \leq, \ll, \between, \Lambda)$ pour
provenir d'un atelier où $\Lambda$ est dense :
\begin{enumerate}
\item[\textbf{M$\Lambda$ 0}] $X \leq Y \Rightarrow X \ll Y$ ;
\item[\textbf{M$\Lambda$ 1}] (raffinements induits) si $Y \ll X$, $X' \leq X$, et
$Z \in \Lambda$ raffine $Y$ et $X'$, il existe $Y' \leq Y$ avec $Y' \ll X'$ et $Z
\ll Y'$\footnote{Nous énonçons M$\Lambda$ 1 avec $X' \leq X$, comme la page 39 et
comme ML 1. À partir de la page 42, et encore pages 51, 52 et 59, la page dispose
$X'$ sous $X$ avec un signe vertical que la transcription lit $\ll$. Sous cette
lecture l'axiome demanderait que les strates d'un raffinement $Y$ qui raffinent un
\emph{autre} raffinement $X'$ couvrent leur ombre commune, ce que les ateliers
polyédraux ne vérifient pas (deux petits triangles de deux subdivisions d'un même
triangle qui se chevauchent : un point intérieur aux deux n'est dans aucune
strate du premier contenue dans le second). La clause « si $Y \leq X$, on peut
prendre $Y' \leq X'$ » et l'égalité $\widetilde{Y} \cap \mathrm{Omb}(X') =
\widetilde{Y} \cap \widetilde{X}'$ affirmée page 59 ne valent, elles, que pour $X'
\leq X$, par At 5. C'est un point à vérifier sur le facsimilé.} ;
\item[\textbf{M$\Lambda$ 2}] (densité) tout $X$ a un raffiné intérieur dans
$\Lambda$ ;
\item[\textbf{M$\Lambda$ 3}] (strates minimales) si $Z \in \Lambda$ raffine $X$, il
existe $X' \leq X$ avec $Z \mathrel{\overset{\circ}{\ll}} X'$ ;
\item[\textbf{M$\Lambda$ 4}] (saturation) $\between$ est réflexive, symétrique et
héréditaire pour $\leq$ ;
\item[\textbf{M$\Lambda$ 5}] si $X \between Y$ et $Z \in \Lambda$ raffine $X$ et
$Y$, il existe $Z'$ avec $Z \ll Z' \leq X, Y$.
\end{enumerate}
Seuls les deux derniers font intervenir $\between$. M$\Lambda$ 5 s'écrit $X
\between Y \Rightarrow \mathcal{R}(X, Y)$, où $\mathcal{R}$ est la relation,
définie par $\leq$ et $\ll$ seuls,
\[
\mathcal{R}(X, Y) \iff \mathrm{Omb}(X) \cap \mathrm{Omb}(Y) = \bigcup_{Z' \leq X,
Y} \mathrm{Omb}(Z') ,
\]
réflexive, symétrique et héréditaire. Les compatibilités admissibles sont donc les
relations saturées contenues dans $\mathcal{R}$, et la plus large est
$\mathcal{R}$ elle-même --- qui est la compatibilité dans un atelier ensembliste
spécieux, mais pas dans d'autres cas intéressants : $\mathcal{R}(x, y)$ est
toujours vraie pour deux lieux, alors que hors du cas ensembliste deux lieux
peuvent être incompatibles\footnote{La Proposition est énoncée, la démonstration
renvoyant aux raisonnements des pages 25 à 33 ; elle n'est pas refaite.}.

Il cherche ensuite (pages 43 et 44) à définir la compatibilité par la disjonction
sur $\Lambda$ (At C$_{\Lambda}$), puis à axiomatiser $(\mathcal{M}, \leq, \ll,
\Lambda, \parallel_{\Lambda})$ ; la page 44 est entièrement barrée, et la page 47
en reprend les axiomes. Il y apparaît une notion de partie « très dense » : toute
multistrate majore un élément de $\Lambda$.

\pagerange{45}{49}
\subsection*{Ateliers $\Lambda$-spécieux (24 juin, pages 45 à 49)}

Le 24 juin, pour $\Lambda$ fermée pour $\leq$, il pose les analogues des axiomes
des lieux :
\begin{enumerate}
\item[\textbf{At$\Lambda$ 1}] (densité) $\mathrm{omb}_{\Lambda}(X)^{\circ} \neq
\varnothing$ ;
\item[\textbf{At$\Lambda$ 2}] $X \parallel Y$ si et seulement si tout élément de
$\mathrm{omb}_{\Lambda}(X)$ est disjoint de tout élément de
$\mathrm{omb}_{\Lambda}(Y)$ ;
\item[\textbf{At$\Lambda$ 3}] (ou At D$\Lambda$) pour $X' \in
\mathrm{omb}_{\Lambda}(X)$, $Y' \in \mathrm{omb}_{\Lambda}(Y)$, $X'_{L} \parallel
Y'_{L} \Rightarrow X' \parallel Y'$ --- c'est At L 3 si $\Lambda = \mathcal{L}$ ;
\item[\textbf{At C$\Lambda$ (spéc)}] $X \between Y$ si et seulement si les éléments
de $\mathrm{omb}_{\Lambda}(X)$ et de $\mathrm{omb}_{\Lambda}(Y)$ qui évitent $L$
sont disjoints ; \textbf{At C$\Lambda$ (pol)} ajoute que $L$ est vide ou une
multistrate.
\end{enumerate}
Avec At$\Lambda$ 1 et At$\Lambda$ 2, $\Lambda$ est dite très dense. Le tableau
d'implications de la page 22 se transpose (page 45), et la page 46 donne quatre
systèmes complets équivalents d'axiomes de compatibilité-disjonction, par
exemple $\{\Lambda 1, \Lambda 2, \Lambda 3, C\Lambda^{*}\}$ ou $\{\Lambda 1,
\Lambda 2, C^{*}, D_{0}\}$, l'astérisque désignant indifféremment la version
spécieuse ou polyédrale ; $\Lambda 1$ est le seul à figurer dans chacun. Quand
ils sont satisfaits l'atelier est dit \emph{$\Lambda$-spécieux}, resp.
\emph{$\Lambda$-polyédral}. Pour $\Lambda = \mathcal{L}$, c'est le cas déjà
traité, avec la divisibilité ; pour $\Lambda = \mathcal{M}$, $\Lambda 1$ est
tautologique, $\Lambda 2$ équivaut à $D_{0}$, $\Lambda 3$ à $D$, $C\Lambda^{*}$ à
$C^{*}$, et « $\mathcal{M}$-spécieux » veut dire spécieux tout court, avec $D$.

Les pages 47 à 49 tentent de reconstruire l'atelier à partir de
$\parallel_{\Lambda}$ seule, en définissant $\between$ et $\parallel$ sur
$\mathcal{M}$ par les éléments de $\Lambda$ qui évitent $L$. Le Lemme 1 de la page
48 passe : $X \parallel Y$ équivaut à $X \between Y$ et $\widetilde{X} \cap
\widetilde{Y} = \varnothing$, et sur $\Lambda$ la disjonction définie est celle de
départ. Le Lemme 2 ($X'_{L} \parallel Y'_{L} \iff X' \parallel Y'$) et le Lemme 3
($\mathrm{omb}_{\Lambda}(X) \cap \mathrm{omb}_{\Lambda}(Y) =
\mathrm{omb}_{\Lambda}(L)$ pour $X \between Y$) ne passent pas : il faudrait les
poser en axiomes, et alors l'énoncé revient « tautologiquement » à la Proposition
de la page 42. « Rien de nouveau ! On tourne en rond ! »\footnote{La page 47
écrivait « je dis que ça suffit » ; une note en marge, « Non » doublement
souligné, renvoie à la page 49.}

\pagerange{50}{53}
\section*{Préateliers (pages 50 à 53)}

La page 50 est une récapitulation. Les axiomes en termes de $(\mathfrak{F}, \leq,
\ll)$ sont ceux de la page 23 ; les implications entre axiomes de
compatibilité-disjonction, celles des pages 22, 45 et 46, avec une partie dense
$\Lambda$ pour « clef ». Les données $(\mathcal{M}, \leq, \ll, \between)$ sont plus
maniables, et il propose de les regarder comme une structure algébrique
indépendante, un \emph{préatelier} (ou « algèbre de multistrates »). Un préatelier
peut correspondre à plusieurs ateliers, réalisés comme des $\mathfrak{F} \subset
\mathfrak{P}(\mathcal{M})$, parmi lesquels un plus grand : toutes les parties de
$\mathcal{M}$ fermées pour $\leq$ et formées d'éléments deux à deux compatibles.
Il en donne trois axiomatiques.

\emph{(1) En termes de $(\mathcal{M}, \leq, \ll, \between)$} : Prat 1 ($\leq
\Rightarrow \ll$), Prat 2 (transitivité des raffinements intérieurs), Prat 3
(tout raffinement est raffinement intérieur d'une strate ; $\between$ réflexive,
symétrique, héréditaire), Prat 4 (deux multistrates compatibles dont un même
élément est raffinement intérieur sont égales), Prat 6 (si $X \between Y$, $X \ll Y \Rightarrow X \leq
Y$)\footnote{Il n'y a pas de Prat 5 : « Prat 4 » est écrit au-dessus d'un « Prat
5 » non nettement biffé, et la ligne sur $\between$ n'a pas d'étiquette propre.
Ce sont At 1 à At 7 traduits, sans plus.}. Seuls les trois derniers parlent de
$\between$.

\emph{(2) Avec une partie dense $\Lambda$} : Prat$\Lambda$ 1 à Prat$\Lambda$ 6, qui
sont M$\Lambda$ 0 à M$\Lambda$ 5 de la page 42\footnote{Prat$\Lambda$ 6 est écrit
« $Z' \ll Z \leq X, Y$ » ; on suit M$\Lambda$ 5 : $Z \ll Z' \leq X, Y$. Les
étiquettes Prat$\Lambda$ 5 et 6 sont en surcharge, et Prat$\Lambda$ 4, écrit plus
bas, est à insérer avant, comme le marque une flèche.}.

\emph{(3) Avec une disjonction sur les lieux} : Prat dis $1'$ à $6'$, qui sont ML
0 à ML 5 de la page 34 dans un autre ordre, la compatibilité étant définie par (1) --- à quoi s'ajoute,
dans le cas polyédral, que $L$ soit vide ou de la forme $\widetilde{Z}$, et
l'axiome \textbf{Prat pol} : chaque $\widetilde{X}$ est polyédral\footnote{La page
écrit « si $\{X, Y\}$ est majoré et minoré, alors $\mathrm{Inf}^{\leq}(Y, Z)$
existe » ; il s'agit de deux éléments $Y, Z$ de $\widetilde{X}$ qui ont un minorant
commun. La borne inférieure existe alors dans $\widetilde{X}$, note-t-il, « mais
pas nécessairement dans $\mathcal{M}$ ».}. « Donc le \emph{vrai} problème
d'axiomatisation », dit une note en marge, est ailleurs --- la suite n'est pas
lue.

\textbf{Remarque} (page 53). Un ensemble ordonné $(\mathcal{M}, \leq)$
quelconque, avec $\ll = \leq$, la compatibilité toujours satisfaite et $x
\parallel y \iff x \neq y$ sur les éléments minimaux, est un préatelier
« trivial » : $X \mathrel{\overset{\circ}{\ll}} Y$ y signifie $X = Y$, les figures
sont toutes les parties fermées, $F \parallel G \iff F \cap G = \varnothing$. Il
satisfait At D, At C (spéc), At Fil, At supp, At L 3, et At L 2 dès que tout
élément majore un élément minimal (ce qui est le cas d'un ensemble ordonné fini)\footnote{La
page dit d'At L 2 qu'il « n'est pas prédisant », lecture incertaine ; elle conclut
que le seul axiome non satisfait est At L 1.} --- tous les axiomes, donc, sauf la
divisibilité. C'est le cas de l'ensemble des simplexes d'un ensemble quasi
simplicial, figures et sous-ensembles quasi simpliciaux se correspondant. La
divisibilité est ce qui sépare un complexe combinatoire d'un espace.

\pagerange{54}{60}
\section*{Ateliers ensemblistes et quasi-ensemblistes (pages 54 à 60)}

Le 25 juin, il revient sur les ateliers \emph{ensemblistes}\footnote{Ceux de la
deuxième mouture (dossier 156-6, p.~52-58 : At ens 1 et ses variantes), désormais
caractérisés par un morphisme vers $\mathrm{Fig}(\mathcal{L})$.}. Pour un atelier
$\mathcal{A} = (\mathfrak{F}, \leq, \ll)$, on a l'homomorphisme canonique
\[
\varphi : \mathfrak{F} \to \mathrm{Fig}(\mathcal{L}), \qquad F \mapsto
\mathrm{multomb}(F) ,
\]
et deux conditions :
\begin{enumerate}
\item[(i)] $\varphi$ est un plongement d'ateliers\footnote{Au sens de la deuxième
mouture (dossier 156-6, p.~51).} d'image un sous-atelier spécieux, c'est-à-dire
a) pour toute figure $F$, $\varphi$ induit une bijection des sous-figures de $F$
sur celles de $\varphi(F)$ ; b) $\varphi(F) \ll \varphi(G) \Rightarrow F \ll G$, d'où
l'injectivité ; c) $\varphi(F) \between \varphi(G) \Rightarrow F \between G$ ;
\item[(ii)] $\mathfrak{F}$ est spécieux (pour $\Lambda = \mathcal{L}$, par exemple
At L 1, At L 2, At L 3, At CL (spéc)) et deux lieux distincts sont toujours
disjoints, de sorte que At L 3 devient automatique\footnote{Le signe de la page est
un $\parallel$ barré de traits horizontaux, que la transcription lit $\nparallel$
(« non disjoints »). La démonstration qui suit établit au contraire $x \neq y
\Rightarrow x \parallel y$, et la définition des ateliers quasi-ensemblistes (page
59) demande $x \parallel y \iff x \neq y$ ; nous lisons donc « disjoints ». C'est à
vérifier sur le facsimilé.}.
\end{enumerate}
Ce que les pages établissent est ceci.
\begin{enumerate}
\item[--] (i) $\Rightarrow$ (ii), et $\varphi$ est injective. Deux lieux distincts ont
pour images des figures ponctuelles distinctes, compatibles, donc ils sont
compatibles par c), donc disjoints. At L 2 : si les ombres de $X$ et de $Y$ ne se
rencontrent pas, $\varphi(X) \parallel \varphi(Y)$, d'où $X \between Y$ par c), et
$X \cap Y = \varnothing$ par At L 1. At CL (spéc) : la condition équivaut à
$\mathrm{omb}(X) \cap \mathrm{omb}(Y) = \mathrm{omb}(L)$, qui fait de $\varphi(L)$
une sous-figure commune et donne $\varphi(X) \between \varphi(Y)$\footnote{Pour At
L 1 la page écrit seulement que $\varphi(X) \neq \varnothing$, « ce qui signifie
$\mathrm{omb}(X)^{\circ} \neq \varnothing$ » ; $\varphi(X) \neq \varnothing$ ne
donne que $\mathrm{omb}(X) \neq \varnothing$, et c'est la bijection a) sur les
sous-figures qui doit exclure que l'ombre de $X$ soit réunion de celles de ses
strates propres. Le pas n'est pas écrit.}.
\item[--] (ii) $\Rightarrow$ (i a). C'est le Lemme 1 de la page 55 : pour $\Lambda$
dense, $\varphi_{\Lambda} : F \mapsto \mathrm{multomb}_{\Lambda}(F)$ induit une
bijection des sous-figures, $\widetilde{F} \to \widetilde{\varphi_{\Lambda}(F)}$
étant un isomorphisme d'ensembles ordonnés --- « vu pour $\Lambda = \mathcal{M}$ ;
dans le cas général, c'est un fait ensembliste ».
\item[--] (ii) et l'injectivité de $\varphi$ sur $\mathcal{M}$ donnent (i c) : la
plus grande sous-figure commune de $\varphi(F)$ et $\varphi(G)$ est alors
$\varphi(L)$, et $\varphi(F) \between \varphi(G)$ équivaut à $\mathrm{omb}(F) \cap
\mathrm{omb}(G) = \mathrm{omb}(L)$, qui est la condition de At CL (spéc).
\item[--] Sans injectivité, (i c) peut manquer. \emph{Exemple} (page 57) : dans un
espace topologique, les simplexes \emph{affines}, avec la relation de face et le
raffinement par inclusion (en comparant les structures affines) ; deux structures
affines différentes sur une même partie homéomorphe à $[0,1]$ donnent $X \neq Y$
de même bord, $\varphi(X) = \varphi(Y)$, et ni $X \between Y$ ni $X \ll Y$.
\item[--] (i b) n'est pas atteint : ramené aux multistrates, « c'est faux ! », et
le contre-exemple de la page 58 prend les figures polyédrales
$\mathrm{Figpol}(\mathcal{L})$ avec le raffinement polyédral et la compatibilité
ensembliste --- deux strates polyédrales avec $X \ll Y$ au sens ensembliste mais
non au sens polyédral\footnote{Le paragraphe qui ouvrait ce contre-exemple (« $X$
espace affine sur $\mathbb{R}$ \ldots ») est barré. La page 56 contient aussi une
démonstration de c) sans injectivité, abandonnée au profit du Corollaire de la page
57 ; dans sa dernière phrase, « mais $X \between Y$ » se lit là où l'argument
attend $X$ et $Y$ non compatibles.}.
\end{enumerate}

\textbf{Définition} (page 59). Un atelier est \emph{ensembliste spécieux} s'il
satisfait (i), \emph{ensembliste polyédral} si l'on remplace
$\mathrm{Fig}(\mathcal{L})$ par $\mathrm{Figpol}(\mathcal{L})$ ; il est
\emph{quasi-ensembliste} s'il satisfait (ii). Les préateliers quasi-ensemblistes
sont donnés par $(\mathcal{M}, \leq, \ll)$ et les quatre axiomes Prat quens 1 à
4 : $\leq \Rightarrow \ll$ ; divisibilité ; strates minimales ; raffinements
induits, $\mathrm{omb}(Y) \cap \mathrm{omb}(X') = \bigcup_{Z \in Y_{X'}}
\mathrm{omb}(Z)$ avec $Y_{X'} = \widetilde{Y} \cap \mathrm{Omb}(X')$. La
disjonction sur les lieux n'a plus à être donnée : c'est $x \neq y$.

La page 60, écrite vite et lue en partie, dit pourquoi cette structure lui
convient : elle est commune aux contextes spécieux et polyédral, elle ne demande
que $\leq$ et $\ll$, et quatre axiomes au lieu de beaucoup plus. Pour aller au-delà
du cas quasi-ensembliste, il faudra une notion de partie fermée « admissible »,
définie par des conditions de finitude locale. Une note en marge, presque
illisible, essaie des noms --- « Préatelier », « Stratologie »,
« stratologique » --- et s'achève sur « topologie ! ».

\pagerange{61}{68}
\section*{Supports (pages 61 à 68)}

\subsection*{La fermeture de Galois de la disjonction}

Soient $\Sigma_{\mathfrak{F}}$ l'ensemble des parties de $\mathfrak{F}$ fermées vers
le bas pour $\ll$ et stables par sommes, et $\Sigma_{\mathcal{M}}$ l'ensemble des
parties de $\mathcal{M}$ fermées vers le bas pour $\ll$. Les applications
$\mathfrak{S} \mapsto \mathfrak{S} \cap \mathcal{M}$ et $S \mapsto \{F \mid
\widetilde{F} \subset S\}$ sont des isomorphismes d'ensembles ordonnés inverses
l'un de l'autre, et l'on identifie les deux. Pour $A \subset \mathfrak{F}$ on
pose
\[
\mathrm{cosupp}\, A = \{G \in \mathfrak{F} \mid G \parallel F \ \text{pour tout}\
F \in A\} , \qquad \mathrm{supp}\, A = \mathrm{cosupp}(\mathrm{cosupp}\, A) ,
\]
qu'il note aussi $\overline{A}$ --- « ne pas confondre avec l'adhérence de $A$ pour
la topologie associée à $\leq$ ou à $\ll$ ! ». Le cosupport est dans
$\Sigma_{\mathfrak{F}}$\footnote{La page invoque un axiome lu « At D » sans
certitude ; At D$_{0}$ suffit pour que $\mathrm{cosupp}\, A$ soit fermé pour
$\ll$, et la stabilité par sommes vient de ce que la disjonction se teste strate
par strate.}. Tout cela est formel : il ne s'agit que d'une relation symétrique
sur un ensemble, et l'on écrit $A \parallel B$ pour $A \subset \mathrm{cosupp}\,
B$, ce qui équivaut à $B \subset \mathrm{cosupp}\, A$ et à $X \parallel Y$ pour
tous $X \in A$, $Y \in B$.

Les propriétés (i) à (x) des pages 62 et 63 sont celles d'une telle fermeture :
$\mathrm{cosupp}$ renverse l'inclusion ; $A \subset \overline{A}$,
$\overline{\overline{A}} = \overline{A}$ ; $\overline{A}$ est le plus petit
\emph{$\mathfrak{F}$-support} --- partie égale à sa fermeture --- contenant $A$ ;
l'ensemble $\Sigma^{\circ}_{\mathfrak{F}}$ des supports est un treillis complet,
d'infimum l'intersection et de supremum la fermeture de la réunion ; $\complement
A = \mathrm{cosupp}\, A$ est un anti-isomorphisme involutif de
$\Sigma^{\circ}_{\mathfrak{F}}$ (parce que $\parallel$ est symétrique), qui échange
le plus petit support $\{\varnothing_{\mathfrak{F}}\}$ et le plus grand
$\mathfrak{F}$ ; enfin $A \cap \complement A = \{\varnothing_{\mathfrak{F}}\}$,
parce que la figure vide est la seule disjointe d'elle-même, et donc $A \subset B
\Rightarrow A \cap \complement B = \{\varnothing_{\mathfrak{F}}\}$\footnote{Ce
sont les axiomes d'un \emph{treillis orthocomplémenté}, et la construction est
celle des fermés d'une polarité (G.~Birkhoff, 1940) ou d'un \emph{espace
d'orthogonalité} (J.~R.~Dacey, 1968) ; la relation de disjonction y joue le rôle
de l'orthogonalité. La première mouture (dossier 156-4, p.~65-68) définissait déjà
cosupport, support et « supports virtuels » de la même façon, mais présumait une
algèbre de Boole ; c'est ici que la structure booléenne est mise en doute. Les
noms sont les nôtres.}.

\emph{NB} (pages 63 et 64). La réciproque de la dernière implication n'est pas vraie
en général (« il est faux »), la réunion de deux supports n'est pas \emph{a priori}
un support --- une « figure » contenue dans la réunion n'est pas forcément contenue
dans l'un ou dans l'autre --- et la distributivité $P \wedge \bigvee_{i} Q_{i} =
\bigvee_{i} (P \wedge Q_{i})$ n'a pas de raison d'être vraie, même pour des sommes
finies. « Pourtant, cela a tendance à être vrai » dans les ateliers auxquels il a
songé. Les pages 81 à 83 montreront le contraire.

\subsection*{Changer d'ensemble de référence}

Tout se redit en termes de la disjonction dans $\mathcal{M}$ : les
$\mathcal{M}$-supports sont en bijection canonique avec les $\mathfrak{F}$-supports,
compatible au complément, et cette fois le plus petit $\mathcal{M}$-support est la
partie vide. Plus généralement, pour une partie $\Lambda \subset \mathfrak{F}$
satisfaisant
\[
\mathrm{Disj}_{\Lambda} : \quad F \parallel G \iff X \parallel G \ \text{pour
tout}\ X \in \mathrm{omb}_{\Lambda}(F) ,
\]
on a $\mathrm{cosupp}\, A = \mathrm{cosupp}\, \mathrm{omb}_{\Lambda}(A)$ pour toute
partie $A \subset \mathfrak{F}$\footnote{La page écrit « $A \in \mathfrak{F}$ ».},
$\mathrm{omb}_{\Lambda}(A) = A \cap \Lambda$ si $A$ est fermée pour $\ll$, et une
bijection $\Sigma^{\circ}_{\mathfrak{F}} \simeq \Sigma^{\circ}_{\Lambda}$, $A \mapsto
A \cap \Lambda$, d'inverse $S \mapsto \{F \mid \mathrm{omb}_{\Lambda}(F) \subset
S\}$, qui respecte l'ordre et le complément. Pour comprendre les supports on a
intérêt à prendre $\Lambda$ petit : $\Lambda = \mathcal{L}$ s'impose, si At L 2 est
satisfait.

Le cas « particulièrement agréable » est celui où la disjonction des lieux est $x
\neq y$, par exemple un atelier quasi-ensembliste. Alors $\Sigma_{\mathcal{L}} =
\mathfrak{P}(\mathcal{L})$, le complément est le complément ordinaire, le support
d'une partie $A$ est $\mathrm{omb}(A)$, la disjonction devient
\[
F \parallel G \iff \mathrm{omb}\, F \cap \mathrm{omb}\, G = \varnothing ,
\]
et les perplexités se dissipent : $A \subset B \iff A \cap \complement B =
\varnothing$, et la distributivité vaut. La page 68 --- barrée d'un trait oblique
--- généralise à une partie $P \subset \mathcal{L}$ de « points », deux à deux
disjoints et qui détectent la disjonction ; les supports s'identifient alors à
$\mathfrak{P}(P)$, de façon compatible au complément, et le support d'une partie
est la réunion des ensembles de points de ses éléments\footnote{La condition b),
lourdement raturée, est énoncée de façon qu'une note en marge juge elle-même « un
peu brève ! ». La page écrit $\mathrm{omb}_{\Lambda}$ où l'on attend
$\mathrm{omb}_{P}$, et « $A$, $B$ » pour $F$, $G$.}.

\pagerange{69}{80}
\section*{Réalisations par des points (pages 69 à 80)}

\subsection*{Un doute, et l'exemple d'un ensemble ordonné (25 juin, pages 69 à 72)}

« En dehors du cas quasi-ensembliste, je suspecte que la notion de support,
telle qu'elle est présentée ici, est caduque » : on ne devrait parler du support
d'une partie $A$ que si ses éléments sont mutuellement compatibles. Exemple
typique : les strates d'un ensemble ordonné $I$. Les 0-strates sont les points,
les 1-strates les segments $[a, b] = I_{a \leq \cdot \leq b}$ avec $a < b$, de bord
$\partial X = \{a, b\}$ ; $X \leq Y$ si $X = Y$ ou si $X$ est une extrémité de $Y$
; $X \ll Y$ si $X \subset Y$ ; et $X \between Y$ dans trois cas exclusifs : $X =
Y$ ; l'un est extrémité de l'autre ; ou $\partial X \cup \partial Y$ est
totalement ordonné et $X \cap Y$ est vide ou réduit à une extrémité commune
(segments séparés ou bout à bout). Tous les axiomes sont satisfaits, dit-il, sauf
At L 1, qui l'est si et seulement si les intervalles ouverts $]a, b[$ ne sont pas
vides --- l'ordre est dense. Les figures finies sont les parties finies, fermées,
d'éléments deux à deux compatibles ; les strates y sont combinatoirement finies,
$\widetilde{X}$ ayant un élément pour un point et trois pour un
segment\footnote{La page porte, pour un point, un chiffre lu $0$, petit et peu
formé ; $\widetilde{X} = \{X\}$ en a un.}.

Au chapitre III, il avait défini le support d'une telle figure comme partie de
$I$, et appelé « modérées fermées » des parties qu'il aurait dû appeler
« modérées », réservant « constructibles » à celles qu'il appelait
« modérées »\footnote{Dossier 156-3 ; « loc.\ cit.\ p.~8 » y renvoie. La phrase
du milieu de la page 70 n'est lue qu'en partie.}. Cela donne un bon formalisme
de réunion, d'intersection et de complément pour des parties constructibles deux
à deux compatibles --- dont la réunion des bords est totalement ordonnée. Sans
hypothèse sur $I$, il faut compter avec des segments $[a, b]$ réduits à $\{a,
b\}$, et il serait déraisonnable de leur donner $[a, b]$ pour support. Mais à une
chaîne finie $a_{1} < \cdots < a_{n}$ on peut associer l'algèbre de Boole engendrée
par les « points » $\{a_{1}\}, [a_{1}, a_{2}], \{a_{2}\}, \ldots, \{a_{n}\}$
--- les cellules ouvertes ---, et définir les supports des familles qui lui sont
subordonnées. Pour une figure $F$, les supports forment $\mathfrak{P}(\widetilde{F})$,
et un raffinement $F' \ll F$, par l'application $\varphi : \widetilde{F}' \to
\widetilde{F}$ qui envoie une strate sur la plus petite strate qu'elle raffine,
induit l'image réciproque $\varphi^{*} : \mathfrak{P}(\widetilde{F}) \to
\mathfrak{P}(\widetilde{F}')$. Ce qu'il ne croit pas, hors du cas
quasi-ensembliste, c'est qu'on puisse interpréter ces $\mathfrak{P}(\widetilde{F})$,
pour des figures sans rapport, comme parties d'un même $\mathfrak{P}(P)$.

\subsection*{Réalisations par supports (pages 72 à 80)}

Il définit pourtant une \emph{réalisation par supports} d'un atelier dans un
ensemble $P$ de « points » : une application $\mathcal{M} \to \mathfrak{P}(P)$, $X
\mapsto |X|$, telle que
\begin{enumerate}
\item[(i)$_{\mathcal{M}}$] $X \ll Y \Rightarrow |X| \subset |Y|$ ;
\item[(ii)$_{\mathcal{M}}$] $X \between Y \Rightarrow |X| \cap |Y| = \bigcup_{Z \leq
X, Y} |Z|$ ;
\item[(iii)$_{\mathcal{M}}$] $X \parallel Y \iff |X| \cap |Y| = \varnothing$ --- « le
fond de l'axiome est dans $\Leftarrow$ » ;
\item[(iv)$_{\mathcal{M}}$] une subdivision ne change pas $|X|$ ;
\item[(v)$_{\mathcal{M}}$] $|X| = \bigcup_{Z \in \mathrm{omb}(X)} |Z|$ ;
\item[(vi)$_{\mathcal{M}}$] (a) les $|X|$ recouvrent $P$, et (b) ils séparent les
points\footnote{La page écrit (b) avant (a), et un trait renvoie (a) en tête.}.
\end{enumerate}
(vi) est une normalisation : on peut toujours passer au quotient de $P$ par la
relation « appartenir aux mêmes $|X|$ ». Par (v) il suffit de connaître $|x|$
pour les lieux, et la variante en termes de lieux demande que $x \parallel y \iff
|x| \cap |y| = \varnothing$ et que les $|x|$ séparent les points. Dans tous les
cas intéressants connus --- sauf peut-être celui d'un ensemble ordonné donné
\emph{ad hoc} ---, on trouve une notion naturelle de support ponctuel ; dans
l'exemple d'une ficelle topologique, la séparation des points peut manquer. La
question est alors : une telle donnée détermine-t-elle $P$ à isomorphisme unique
près, par un isomorphisme $\Sigma_{\mathcal{M}} \simeq \mathfrak{P}(P)$ compatible
au complément ?

Pour $A \subset \mathcal{M}$ on pose $|A| = \bigcup_{X \in A} |X|$. Avec (iii) seul,
$X \in \mathrm{cosupp}\, A$ équivaut à $|X| \cap |A| = \varnothing$, d'où :

\textbf{Lemme} (page 76). Sous (iii)$_{\mathcal{M}}$, $|A| \subset |B| \Rightarrow
\mathrm{supp}\, A \subset \mathrm{supp}\, B$ ; la réciproque vaut si $B$ est un
support. \textbf{Corollaire.} $A \mapsto |A|$ est injective sur les supports et y
induit l'inclusion.

En général $|\bigwedge A_{i}| \subset \bigcap |A_{i}|$, $|\bigvee A_{i}| \supset
\bigcup |A_{i}|$ et $|\complement A| \subset P \smallsetminus |A|$, sans égalité
garantie, et l'on n'a que $|A| \subset |\mathrm{supp}\, A|$. « Tout se résout »
sous l'hypothèse

\textbf{Hyp.} Si $x \notin |A|$, il existe $X \parallel A$ avec $x \in |X|$ ---
c'est-à-dire $|\mathrm{cosupp}\, A| = P \smallsetminus |A|$ pour toute partie $A$.

\textbf{Lemme 2} (page 78). Sous (iii)$_{\mathcal{M}}$, (vi a)$_{\mathcal{M}}$ et
Hyp., on a $|A| = |\mathrm{supp}\, A|$ pour toute partie $A$, et
$\mathrm{supp}\, A \subset \mathrm{supp}\, B \iff |A| \subset |B|$ ; l'injection
$\Sigma_{\mathcal{M}} \to \mathfrak{P}(P)$ commute aux compléments\footnote{En
effet $|\mathrm{supp}\, A| = |\mathrm{cosupp}(\mathrm{cosupp}\, A)| = P
\smallsetminus |\mathrm{cosupp}\, A| = P \smallsetminus (P \smallsetminus |A|) =
|A|$ ; la page écrit « $= A$ », et « $\overline{X}$ » pour $|X|$ à la page 77.}.

Pour les intersections il faut plus : \textbf{Hyp'}, chaque point est le support
$|X|$ d'une multistrate. Elle entraîne Hyp. et la surjectivité, et donne
$|\bigcap A_{i}| = \bigcap |A_{i}|$ pour des supports : si $|X| = \{x\}$ est contenu
dans chaque $|A_{i}|$, alors $X \in A_{i}$ pour tout $i$.

\textbf{Proposition} (page 80). Soit $X \mapsto |X|$ une application $\mathcal{M}
\to \mathfrak{P}(P)$ telle que
\[
\mathrm{Supp}^{1}_{P} : X \parallel Y \iff |X| \cap |Y| = \varnothing, \qquad
\mathrm{Supp}^{2}_{P} : \text{tout point est un } |X| .
\]
Alors $|A| = |\mathrm{supp}\, A|$ et $|\mathrm{cosupp}\, A| = P \smallsetminus |A|$
pour toute partie $A$\footnote{La page écrit $|A| = \mathrm{supp}\, A$, sans les
barres au second membre ; c'est l'égalité du Lemme 2.} ; $\mathrm{supp}\, A
\subset \mathrm{supp}\, B \iff |A| \subset |B|$, et $A \parallel B \iff |A| \cap
|B| = \varnothing$ --- ceci sans $\mathrm{Supp}^{2}_{P}$ ; enfin $A \mapsto |A|$ est
un isomorphisme d'ensembles ordonnés $\Sigma_{\mathcal{M}} \simeq \mathfrak{P}(P)$
qui commute aux compléments. Les supports forment alors une algèbre de Boole
complète atomique, dont les atomes sont les points\footnote{Le nom est le nôtre.
La surjectivité, que la page n'écrit pas, vient de $\mathrm{Supp}^{2}_{P}$ : pour
$S \subset P$, la famille des $X$ avec $|X| = \{x\}$, $x \in S$, a pour support une
partie d'image $S$. La page 81 ajoute que cela détermine $P$ et $A \mapsto |A|$ à
isomorphisme canonique près.}.

\pagerange{81}{83}
\section*{La ficelle à deux brins (pages 81 à 83)}

La page 81, en prose rapide et lue par fragments, prend pour lieux les parties
compactes contractiles d'un espace $P$, construit des « quasi-simplexes » de
basse dimension par des conditions \emph{ad hoc} et constate que « tout marche ! »
--- sauf que, si l'on regarde les lieux sur une ficelle donnée,
$\mathrm{Supp}^{1}_{P}$ cesse d'être vrai en général.

\emph{Exemple.} La « ficelle à deux brins » : deux arcs joignant un point $0$ à
un point $1$. Un lieu qui la coupe en deux est un couple de points, un sur chaque
brin, et les lieux intérieurs s'identifient à $]0, 1[ \times ]0, 1[$. Plus
généralement, soient $I$, $I'$ deux ensembles totalement ordonnés, et $J = I
\times I'$ ; pour $z = (x, x')$, le support ponctuel est $|z| = \{x, x'\}$, et
$|0| = \{0\}$, $|1| = \{1\}$. Deux lieux sont disjoints lorsqu'ils sont
emboîtés\footnote{C'est notre lecture du modèle, que la page ne décrit pas en
entier ; elle est celle des lieux d'une faille au chapitre II (dossier 156-2,
p.~8), où deux lieux sont « disjoints » s'ils se suivent de $A$ vers $B$, ce qui
est plus fort que l'intersection vide. Elle rend compte de toutes les formules
des pages 82 et 83.} : $z' < z$ ou $z' > z$ dans l'ordre produit. Le support de
$\{z\}$ est $A_{z} = \{z\}$, son complément $B_{z} = \{z' \mid z' < z \text{ ou }
z' > z\}$, et $|B_{z}| = (I \sqcup I') \smallsetminus |z|$.

Mais si $z_{1} = (x_{1}, x'_{1})$ et $z_{2} = (x_{2}, x'_{2})$ se croisent ---
$x_{1} < x_{2}$ et $x'_{1} > x'_{2}$ ---, un lieu comparable aux deux est sous les
deux ou au-dessus des deux, et
\[
|B_{z_{1}} \cap B_{z_{2}}| \subsetneqq |B_{z_{1}}| \cap |B_{z_{2}}| ,
\]
la différence étant $]x_{1}, x_{2}[$ sur $I$ et $]x'_{2}, x'_{1}[$ sur $I'$. Par
complément, pour $A = \{z_{1}, z_{2}\}$,
\[
|A| = \{x_{1}, x_{2}, x'_{1}, x'_{2}\} \subsetneqq |\mathrm{supp}\, A| = [x_{1},
x_{2}] \cup [x'_{2}, x'_{1}] .
\]
Enfin, pour $z_{3} = (x_{3}, x'_{3})$ avec $x_{1} < x_{3} < x_{2}$ et $x'_{2} <
x'_{3} < x'_{1}$, on a $A_{z_{3}} \cap (A_{z_{1}} \vee A_{z_{2}}) = A_{z_{3}}$ tandis
que $(A_{z_{3}} \cap A_{z_{1}}) \vee (A_{z_{3}} \cap A_{z_{2}})$ est le plus petit
support : « on n'a pas la propriété modulaire » ni la distributivité de $\wedge$
par rapport à $\vee$\footnote{Le texte écrit « la propriété modulaire », suivi
d'un mot non lu, à côté de la distributivité. L'exemple réfute la distributivité ;
nous n'affirmons rien de la modularité.}.

\textbf{Conclusion-Scholie.} Le formalisme des supports est bon pourvu qu'on ne
s'en serve pas pour comparer les supports de strates ou de figures qui ne sont
pas compatibles\footnote{Les dernières lignes de la page 83 ne sont lues que par
fragments.}. « C'est cette intuition que je voudrais maintenant essayer de cerner
de plus près. »

\pagerange{84}{100}
\section*{L'intérieur d'une multistrate (pages 84 à 100)}

\subsection*{Strates ouvertes (pages 84 à 87)}

Pour une figure $F$, on voudrait identifier les éléments de $\widetilde{F}$ à des
« strates ouvertes », et $\mathfrak{P}(\widetilde{F})$ à un ensemble de supports,
réunions de strates ouvertes. Il faut pour cela donner un sens intrinsèque à
« $X \smallsetminus \partial X$ ». Première idée :
\[
X^{\circ}_{\mathrm{s}} = \mathrm{supp}\, X \cap \mathrm{cosupp}\, \partial X
\subset \mathcal{M} ,
\]
avec le souhait que $\mathrm{supp}\, X$ soit la borne supérieure, dans les
supports, des $Y^{\circ}_{\mathrm{s}}$, $Y \in \widetilde{X}$ ; c'est-à-dire que
$Z \parallel X$ si et seulement si $Z$ est disjoint de tous les
$Y^{\circ}_{\mathrm{s}}$ --- le sens direct est clair. Moyennant At L 2 et At L 3,
on a
\[
(\delta) \qquad \mathrm{supp}\, \mathrm{omb}(X)^{\circ} \subset
X^{\circ}_{\mathrm{s}} .
\]
\footnote{L'étiquette ($\delta$) est lue sans certitude.}
Il introduit alors \textbf{At supp 1} : si $Y \in X^{\circ}_{\mathrm{s}}$ --- $Y$
disjoint de toute strate propre de $X$, et disjoint de tout ce qui est disjoint
de $X$ --- alors $Y \ll X$. D'où le \emph{Corollaire} (page 86) : moyennant At D,
\[
X^{\circ}_{\mathrm{s}} = \{Y \in \mathrm{Omb}(X) \mid Y_{\partial X} = \varnothing\}
.
\]
Une note en marge datée du 26 juin juge l'axiome « un peu artificiel » : il n'est
pas satisfait par les « gros ateliers », ni par l'atelier linéaire par morceaux.
La \emph{Remarque} de la page 87 le renforce en At supp 1 (spéc)\footnote{Dont
l'hypothèse (i) est écrite « $\mathrm{supp}\, Y \subset X$ », où l'on attend
$\mathrm{supp}\, Y \subset \mathrm{supp}\, X$.}, et suggère qu'il doit exister une
définition raisonnable de $\ll$ au moyen de $\leq$ et de la disjonction --- mais
d'une disjonction entre des objets « virtuels » comme $Y \smallsetminus L$, que
At C (spéc) fait déjà intervenir implicitement.

\subsection*{L'intérieur défini par la grande ombre (pages 88 à 91)}

« Faisons abstraction pour le moment des axiomes de supports. » La
\textbf{Définition} de la page 89 est
\[
X^{\circ} = \{Y \in \mathrm{Omb}(X) \mid Y \parallel \partial X\} ,
\]
ce qui, sous At D, équivaut à $Y_{\partial X} = \varnothing$ : aucune strate de $Y$
ne raffine une strate propre de $X$. Une note avertit que $X^{\circ} \cap Y^{\circ}
= \varnothing$ n'entraîne $X^{\circ} \parallel Y^{\circ}$ que si deux lieux
distincts sont disjoints.

\textbf{Lemme 1} (page 89). $X \parallel Y$ si et seulement si $X'^{\circ}
\parallel Y'^{\circ}$ pour toutes strates $X'$ de $X$ et $Y'$ de $Y$. Le sens
direct est clair par At D$_{0}$ ; la réciproque utilise At L 2 et At L 3\footnote{Le
lemme porte « (?) ». La réciproque se fait ainsi : un lieu $x$ de $X$ est raffiné
intérieur d'une strate $X_{1}$ (At 6), il n'en raffine aucune strate propre, donc
$x_{\partial X_{1}} = \varnothing$ et $x \in X_{1}^{\circ}$ ; de même pour $y$ ; donc
$x \parallel y$, et At L 2 conclut.}.

\textbf{Lemme 2} (page 89). Moyennant At L 1 à At L 3, pour $Y \ll X$ : a) $Y
\mathrel{\overset{\circ}{\ll}} X$ si et seulement si $Y^{\circ} \subset X^{\circ}$ ;
b) $Y \in X^{\circ}$ si et seulement si $Y'^{\circ} \subset X^{\circ}$ pour toute
strate $Y'$ de $Y$. La démonstration de a) utilise, dans un sens, que
$Y_{\partial X} \leq \partial Y$ quand $Y \mathrel{\overset{\circ}{\ll}} X$, et dans
l'autre que $Y^{\circ} \neq \varnothing$, qui vient de la divisibilité ; b) vient de
ce que $X^{\circ}$ est fermé pour $\ll$. Une note dit que $Y \in X^{\circ}$, qu'il
écrit $Y \overset{\circ\circ}{\ll} X$, signifie que l'application $\widetilde{Y} \to
\widetilde{X}$ est constante de valeur $X$.

\emph{NB} (page 91). Si $X \between Y$ et $X \neq Y$, alors $X^{\circ} \cap
Y^{\circ} = \varnothing$\footnote{La page écrit seulement $X \neq Y$, au-dessus
d'une condition surchargée qui pourrait être $X \between Y$. L'hypothèse de
compatibilité est nécessaire : la démonstration utilise $\mathrm{Omb}(X) \cap
\mathrm{Omb}(Y) = \mathrm{Omb}(L)$, et deux structures affines différentes sur un
même segment (page 57) ont des intérieurs communs. Le Lemme 4 de la page 95 pose
l'hypothèse ; une note en marge renvoie à lui : « Mieux, $X^{\circ} \parallel
Y^{\circ}$ ».}.

\subsection*{Un théorème, et le lemme qui devait le donner (pages 91 à 98)}

Soit $\Phi = \widetilde{F}$ une partie finie de $\mathcal{M}$, fermée pour $\leq$,
formée de multistrates deux à deux compatibles. Pour $A \subset \widetilde{F}$ on
pose
\[
S_{A} = \mathrm{Supp}_{\mathcal{M}}\Bigl(\bigcup_{X \in A} X^{\circ}\Bigr) =
\mathop{\mathrm{Sup}}_{X \in A} \mathrm{Supp}(X^{\circ}) ,
\]
d'où une application croissante $\mathfrak{P}(\widetilde{F}) \to
\Sigma_{\mathcal{M}}$. Si l'on admet At supp 1, $X^{\circ}$ est déjà un support.

\textbf{Théorème} (page 93), souhaité : (a) $A \mapsto S_{A}$ est un isomorphisme
de $\mathfrak{P}(\widetilde{F})$ sur une partie ordonnée de $\Sigma_{\mathcal{M}}$ ;
(b) elle commute aux bornes supérieures (trivial) et aux bornes inférieures
quelconques ; (c) pour $B \subset A$, $S_{A \smallsetminus B} = S_{A} \cap
\complement S_{B}$ ; (d) si $A = \widetilde{G}$ est fermée, $S_{A} =
\mathrm{supp}\, G$.

Ce qui est démontré : (d), par le Lemme 1 de la page 89 ; et (a), par le
\emph{Lemme 3} (si $B \cup \{X\}$ est formé d'éléments mutuellement compatibles,
$X^{\circ} \subset S_{B} \Rightarrow X \in B$) et le \emph{Lemme 4} (pour $X
\between Y$ distincts, $X^{\circ} \parallel Y^{\circ}$) : on prend $Z \in
X^{\circ}$, qui existe par la divisibilité, et $Z$ est disjoint de tous les
$Y^{\circ}$, $Y \in B$, mais pas de $X^{\circ}$\footnote{Dans la démonstration du
Lemme 4, la page écrit $Y'_{\partial X}$ où l'on attend $Y'_{\partial Y}$.}.

Pour (b) et (c), il voit le théorème comme cas particulier d'un énoncé formel :
dans un ensemble muni d'une relation symétrique antiréflexive $\parallel$, avec
ses supports, leurs bornes et le complément $\complement$, qui vérifie $S \cap
\complement S = \varnothing$ et $S \vee \complement S = \mathcal{M}$, soit une
famille $(A_{i})$ de parties non vides deux à deux disjointes, et $S_{J}$ la borne
supérieure des $\mathrm{supp}\, A_{i}$, $i \in J$ ; alors $J \mapsto S_{J}$ serait
un isomorphisme sur son image commutant à toutes les bornes, avec $S_{A
\smallsetminus B} = S_{A} \cap \complement S_{B}$. L'injectivité passe (Lemmes 1 et 2
de la page 97). La formule de différence se ramène au \emph{Lemme 3} de la page 97 :
\[
S \parallel S' \Longrightarrow S' = (S \vee S') \cap \complement S .
\]
« Mais c'est faux ! » Dans un graphe où une chaîne $A' - A - B - B'$ est telle que
$A$ et $B$ ne sont liés qu'entre eux et à $A'$, $B'$, on a $\mathrm{cosupp}\, A =
\{A', B\}$, $\mathrm{supp}\, A = \{A\}$, $\mathrm{supp}\, B = \{B\}$\footnote{La page
écrit « $\mathrm{supp}\, A = \mathrm{cosupp}\{A', B\} = A'$ » ; les voisins communs
de $A'$ et de $B$ se réduisent à $A$. Elle écrit aussi « $A \cup B = \complement A
\cap \complement B$ » pour $\complement(A \vee B)$.} ; aucun sommet n'est lié à la
fois à $A$ et à $B$, donc $A \vee B = \mathcal{M}$, et le lemme affirmerait
$\{B\} = \complement A = \{A', B\}$. « Rideau… »\footnote{Le Lemme 3 est la
\emph{loi orthomodulaire} (K.~Husimi, 1937) : pour $S' \leq \complement S$, $S' =
(S \vee S') \wedge \complement S$. Le treillis des fermés d'un graphe pour la
fermeture « voisins des voisins » est toujours orthocomplémenté, rarement
orthomodulaire, et le graphe de la page 98 en est un exemple. Le nom est le nôtre ;
rien n'indique que Grothendieck ait connu ces travaux, issus de la logique
quantique de Birkhoff et von Neumann (1936).}

Il pense pourtant que la formule devrait être vraie, au moins sous la
divisibilité At L 1, « avec l'approche que j'ai prise pour donner un sens
intrinsèque aux “supports constructibles”, via les $X^{\circ}$ ».

\pagerange{99}{100}
\subsection*{L'intérieur défini par les lieux (26 juin, pages 99 et 100)}

Puisqu'il est obligé d'admettre At L 1 (en plus de At L 2 et At L 3), autant
définir plus simplement
\[
X^{\circ}_{\ell} = \mathrm{Supp}\, \mathrm{omb}(X)^{\circ} ,
\]
« c'est pareil » grâce au Corollaire de la page 86 et au fait que, sous At L 2,
tout support $S$ est le support de ses lieux, $S = \mathrm{supp}(S \cap
\mathcal{L})$\footnote{L'identité de $X^{\circ}_{\ell}$ avec $X^{\circ}$ et
$X^{\circ}_{\mathrm{s}}$ passe par At supp 1, que la note du même jour, page 86,
venait de juger artificiel ; sans lui on n'a que l'inclusion ($\delta$).}.

\textbf{Lemme} (page 99). Si $F \preccurlyeq X$ est une subdivision, et
$\widetilde{F}^{\circ} = \{X' \in \widetilde{F} \mid X'
\mathrel{\overset{\circ}{\ll}} X\}$, alors $X^{\circ}_{\ell}$ est la borne
supérieure des $X'^{\circ}_{\ell}$, $X' \in \widetilde{F}^{\circ}$ --- alors même
que, dans le cas quasi-ensembliste, $\mathrm{omb}(X)^{\circ}$ n'est pas la réunion
des $\mathrm{omb}(X')^{\circ}$. L'inclusion $\supset$ est claire, et la réciproque
équivaut au \textbf{Corollaire} (page 100) : si $Z$ est disjoint de tous les
$\mathrm{omb}(X')^{\circ}$, il l'est de $\mathrm{omb}(X)^{\circ}$ ; par At L 2 on se
ramène au cas où $Z$ est un lieu.

Il le vérifie sur une réalisation par points satisfaisant $\mathrm{Supp}^{1}_{P}$
et $\mathrm{Supp}^{2}_{P}$, où $F \preccurlyeq X$ signifie $F \ll X$ et $|F| = |X|$,
et annonce (a) $|X^{\circ}_{\ell}| = |X| \smallsetminus |\partial X|$ et (b)
$\widetilde{F}^{\circ}$ est l'ensemble des $X' \in \widetilde{F}$ dont l'intérieur
est contenu dans celui de $X$, ou le rencontre. La page démontre l'inclusion
$|X^{\circ}_{\ell}| \subset |X| \smallsetminus |\partial X|$ et s'arrête là.

\pagerange{101}{113}
\section*{Magasins (pages 101 à 113)}

\subsection*{Une relation nouvelle (pages 101 à 105)}

« J'ai envie de prendre sur un “magasin” $(\mathcal{M}, \leq, \ll)$ comme
relations primitives, en plus de $\leq$ et $\ll$, et au lieu de $\between$ ou de
$\parallel$ », la relation \emph{intérieurement disjoints}
\[
X \mathrel{|{\circ}|} Y \quad \text{(heuristiquement : } X^{\circ} \parallel Y^{\circ}) ,
\]
de façon à donner un sens à l'intérieur même là où les définitions précédentes le
font vide. L'exemple est celui de la page 53 : pour une maquette simpliciale, avec
$\ll = \leq$, la compatibilité toujours vraie et $X \parallel Y$ si $X$ et $Y$
n'ont pas de face commune, les supports s'identifient aux parties quelconques de
l'ensemble des sommets, le support d'un simplexe étant le simplexe
lui-même\footnote{La page écrit $\Sigma_{\mathcal{M}} \simeq \mathfrak{P}(S)$ sans
dire ce qu'est $S$ ; la disjonction étant l'absence de sommet commun, c'est
l'ensemble des sommets.}. Avec la définition $\mathrm{supp}\, X \cap
\mathrm{cosupp}\, \partial X$, tout simplexe de dimension $\geq 1$ a un intérieur
vide, chacun de ses sommets étant dans une face propre. « C'est triste mais c'est
vrai. »

Avec $\leq$, $\ll$, $\mathrel{|{\circ}|}$ pour relations primitives, on définit,
dans la version qu'il appelle « magasin spécial »\footnote{« Spécial » ici et
page 111, là où les pages précédentes disaient « spécieux » ; la version
« magasin polyédral » est annoncée (« on reviendra bientôt ») et ne vient pas
dans ce dossier.}, pour $L = \widetilde{X} \cap \widetilde{Y}$ la plus grande partie
fermée commune,
\[
X \between Y \iff X' \mathrel{|{\circ}|} Y' \ \text{pour tous}\ X' \in \widetilde{X}
\smallsetminus L,\ Y' \in \widetilde{Y} \smallsetminus L ,
\]
définition qui ne fait pas intervenir $\ll$. Il reprend les axiomes de la page 51
avec $\Lambda = \mathcal{M}$ et en ajoute deux :
\begin{enumerate}
\item[\textbf{Mag 4''}] deux strates distinctes d'une même multistrate sont
intérieurement disjointes ;
\item[\textbf{Mag 5''}] $\mathrel{|{\circ}|}$ est symétrique et antiréflexive, et
stable par raffinement intérieur : $X \mathrel{|{\circ}|} Y$, $X'
\mathrel{\overset{\circ}{\ll}} X$, $Y' \mathrel{\overset{\circ}{\ll}} Y \Rightarrow
X' \mathrel{|{\circ}|} Y'$\footnote{La page conclut « implique $X
\mathrel{|{\circ}|} Y$ », sans les primes.}.
\end{enumerate}
Avec Mag 4'' la compatibilité ainsi définie est héréditaire, et deux multistrates
compatibles distinctes sont intérieurement disjointes. On en déduit Mag 6 : si $X
\between Y$ et $Z \ll X, Y$, il y a $Z'$ avec $Z \ll Z' \leq X, Y$. En effet, soient
$X' \leq X$, $Y' \leq Y$ les strates dont $Z$ est raffiné intérieur ; elles sont
compatibles ; si elles étaient distinctes elles seraient intérieurement disjointes,
et Mag 5'' donnerait $Z \mathrel{|{\circ}|} Z$, contre l'antiréflexivité ; donc $X'
= Y'$ est le $Z'$ cherché\footnote{La page invoque « Mag $\mathcal{M}$5 » pour
obtenir $X' \between Y'$ et le corollaire de Mag 4'' pour conclure $X' = Y'$ ;
l'étape qui conclut est l'antiréflexivité de Mag 5''.}.

On obtient un atelier, dont la disjonction vérifie (\emph{Proposition 1}) : $X
\parallel Y$ si et seulement si toutes les strates de $X$ sont intérieurement
disjointes de toutes celles de $Y$ --- par définition, l'antiréflexivité excluant
une strate commune. La \emph{Proposition 2} (?) examine At C (spéc) : le sens
$\Rightarrow$ est démontré (pages 103 et 104), par la plus petite strate et Mag
5'' ; le sens $\Leftarrow$ demanderait « un axiome supplémentaire genre At L 1 »\footnote{Dans
la démonstration, la page écrit une fois « $X'_{L} \neq \varnothing$ » pour
$X'_{L} = \varnothing$.}. Un \emph{Lemme} montre que sous cette condition C, les
lieux intérieurs de deux raffinements évitant $L$ sont disjoints. D'où la
\emph{Proposition} (?) de la page 104 :
\[
X \mathrel{|{\circ}|} Y \iff x \parallel y \ \text{pour tous}\ x \in
\mathrm{omb}(X)^{\circ},\ y \in \mathrm{omb}(Y)^{\circ} ,
\]
dont le sens $\Rightarrow$ vient de Mag 5'' et dont la réciproque reste ouverte.
Inversement, dans un magasin donné par $\leq, \ll, \between$ et satisfaisant At L 1 à
At L 3, on peut \emph{définir} $\mathrel{|{\circ}|}$ par cette formule ; Mag 4''
vient de At L 3, Mag 5'' de ce que $X' \mathrel{\overset{\circ}{\ll}} X$ entraîne
$\mathrm{omb}(X')^{\circ} \subset \mathrm{omb}(X)^{\circ}$, et la compatibilité
$\between'$ que redéfinit $\mathrel{|{\circ}|}$ coïncide avec $\between$ si et
seulement si At C (spéc) --- équivalent sous At L 2 à At CL (spéc) --- est vrai,
parce que $\mathrm{omb}(X) \smallsetminus \mathrm{omb}(L)$ est la réunion des
$\mathrm{omb}(X')^{\circ}$, $X' \in \widetilde{X} \smallsetminus L$.

\subsection*{Le Scholie, et le changement d'optique (pages 106 à 108)}

\textbf{Scholie.} Un $\between$-magasin satisfaisant At L 1 à At L 3 définit un
$\mathrel{|{\circ}|}$-magasin par la formule ci-dessus. Un
$\mathrel{|{\circ}|}$-magasin définit un $\between$-magasin, de disjonction
« toutes les strates intérieurement disjointes ». Pour que $\mathrel{|{\circ}|}$
provienne d'une relation sur les lieux, il faut la condition
\[
(\ast\ast\ast) \qquad X \mathrel{|{\circ}|} Y \iff x \mathrel{|{\circ}|} y \ \text{pour
tous}\ x \in \mathrm{omb}(X)^{\circ},\ y \in \mathrm{omb}(Y)^{\circ} ,
\]
\footnote{La page l'étiquette « Mag 5'' », nom déjà pris page 103, puis parle de
« Mag 6'' » ; le chiffre est mal formé. On l'appelle $(\ast\ast\ast)$, comme la
formule.}
du type d'un axiome de divisibilité : elle entraîne At L 1 --- si
$\mathrm{omb}(X)^{\circ}$ était vide, on aurait $X \mathrel{|{\circ}|} X$ ---, et
aussi At L 2 et At L 3, et At C (spéc) « semble automatique ». Il conclut à une
« équivalence de catégories » entre $\between$-magasins satisfaisant At L 1 à At L 3
et At C (spéc), et $\mathrel{|{\circ}|}$-magasins satisfaisant
$(\ast\ast\ast)$\footnote{Aucune notion de morphisme n'est définie dans ce dossier
; ce qui est affirmé est une correspondance biunivoque entre les deux espèces de
structure sur un même $(\mathcal{M}, \leq, \ll)$. Les homomorphismes de magasins
viennent dans la quatrième mouture (dossier 156-8, p.~54 et suivantes). Le
diagramme de la page 107 ne trace qu'une flèche, vers le bas.}. Les deux notions
apparaissent comme deux généralisations d'une même structure, trop restrictive par
la seule présence de At L 1.

« On dispose donc de deux généralisations possibles, en direction d'une topologie
sans axiome de divisibilité. Mais c'est la direction $\mathrel{|{\circ}|}$ qui me
semble la plus riche » : elle a toujours un $\between$-magasin associé, et en
\emph{plus} une notion de disjonction intérieure qui exprime une intuition
géométrique que la seule compatibilité ne peut dire en l'absence de At L 1. Dans
ce cadre At C (spéc) paraît artificiel. « Donc finalement, je change d'optique
\ldots{} Je reprends tout, une nouvelle fois ! »\footnote{La page 108, écrite sans
formules, n'est lue qu'en partie.}

\subsection*{Magasins et ateliers spéciaux (pages 109 à 113)}

Un \emph{magasin} (de multistrates) est un ensemble $\mathcal{M}$ muni de trois
relations $\leq$, $\ll$, $\mathrel{|{\circ}|}$, cette dernière se lisant « $X$ et
$Y$ sont intérieurement disjoints » et pouvant s'écrire plus tard $X^{\circ}
\parallel Y^{\circ}$. Les axiomes, après un premier jet barré à la page 109, sont :
\begin{enumerate}
\item[\textbf{Mag 1}] $\leq$ et $\ll$ sont des relations d'ordre et $X \leq Y
\Rightarrow X \ll Y$ ;
\item[\textbf{Mag 2}] si $Z \ll X$, l'ensemble $\{X' \in \widetilde{X} \mid Z \ll
X'\}$ a un plus petit élément ; si c'est $X$, on écrit $Z
\mathrel{\overset{\circ}{\ll}} X$ ;
\item[\textbf{Mag 3}] $\mathrel{|{\circ}|}$ est symétrique, antiréflexive, et stable
par raffinement intérieur ;
\item[\textbf{Mag 4}] deux strates distinctes d'une même multistrate sont
intérieurement disjointes\footnote{Mag 3 et Mag 4 sont Mag 5'' et Mag 4'' des pages
102-103.}.
\end{enumerate}
On \emph{définit} alors
\[
X \between Y \iff X' \mathrel{|{\circ}|} Y' \ \text{pour tous}\ X' \in \widetilde{X}
\smallsetminus \widetilde{X} \cap \widetilde{Y},\ Y' \in \widetilde{Y}
\smallsetminus \widetilde{X} \cap \widetilde{Y} ,
\]
\[
X \parallel Y \iff X' \mathrel{|{\circ}|} Y' \ \text{pour tous}\ X' \in
\widetilde{X},\ Y' \in \widetilde{Y} ,
\]
et $X \parallel Y$ si et seulement si $X \between Y$ et $\{X, Y\}$ n'est pas minoré.
Une \emph{figure} du magasin est une partie fermée pour $\leq$ dont les éléments
sont deux à deux compatibles ; les $\widetilde{X}$ en sont.

Un \emph{atelier spécial} est la donnée $(\mathfrak{F}, \leq, \ll ;
\mathrel{|{\circ}|})$ : At 1 (sommes des familles majorées), At 2 (toute figure est
la somme de ses strates, qui forment le magasin $\mathcal{M}$), At 4 (critère de
raffinement par strates), At 3 (spéc) --- (a) $F \vee G$ existe si et seulement si
les strates de $F$ et de $G$ hors de leur partie commune sont intérieurement
disjointes, (b) $\mathrel{|{\circ}|}$ ne relie que des multistrates ---, plus Mag 1
à Mag 4 sur $(\mathcal{M}, \leq, \ll, \mathrel{|{\circ}|})$. De façon équivalente,
c'est un ensemble $\mathfrak{F} \subset \mathfrak{P}(\mathcal{M})$ de parties
fermées pour $\leq$ (At $\mathcal{M}$ 1), stable par passage aux parties fermées
(At $\mathcal{M}$ 2), qui recouvre $\mathcal{M}$ et contient la partie vide (At
$\mathcal{M}$ 3), et où $\Phi \cup \Psi \in \mathfrak{F}$ si et seulement si $X
\between Y$ pour tous $X \in \Phi$, $Y \in \Psi$ (At $\mathcal{M}$ 4)\footnote{La
page écrit la condition de At $\mathcal{M}$ 4 avec $X$, $Y$ sans quantificateur,
sous une première version barrée.}. Il est « à figures finies » si l'on prend pour
$\mathfrak{F}$ les parties finies fermées d'éléments deux à deux compatibles, ce qui
demande \textbf{Mag fin} : chaque $\widetilde{X}$ est fini.

Le dossier s'arrête au premier tiers de la page 113, sur une remarque : au cours
du travail, c'est « finalement l'atelier qui est au centre de la structure,
alors que l'ensemble $\mathfrak{F}$ de toutes les figures joue un rôle un tantinet
accessoire »\footnote{Les mots « c'est finalement l'atelier » sont ajoutés au-dessus
de la ligne, le premier lu avec doute. Ce qui précède --- les pages 109 à 112, où
tout est porté par $(\mathcal{M}, \leq, \ll, \mathrel{|{\circ}|})$ --- ferait
attendre le magasin ; on garde son mot.}. La quatrième mouture (dossier 156-8)
s'ouvre sur un tableau des types de magasins dont le premier, les « magasins
généraux », est celui-ci : $\mathcal{M}$, $\leq$, $\ll$, $\mathrel{|{\circ}|}$, Mag
1 à Mag 4.

\end{document}
